\chapter{Overcoming the limitations: the F8 program}
\label{ch:overcoming}

\begin{quote}\itshape
This chapter in brief: the limitations declared throughout this document are not footnotes
--- they are the frontier. Each becomes a research line with a discovery hypothesis (the
full ten-line program lives in the project plan). This chapter reports the lines executed
so far, starting with the one that unlocks three others: replacing the linearized generator
with a differentiable neural twin.
\end{quote}

\keyidea{A differentiable neural twin (linearization plus a learned residual) gives full-physics fidelity at microseconds with exact gradients, unlocking the nonlinear fleet --- yet bolting physics features onto a learner fails a second independent way, hardening the warning that ``physics-informed'' is not a free win.}


\section{The program in one view}
\label{sec:f8-program}

Ten declared limitations, ten research lines. The ordering is by synergy, not severity: a
\textbf{differentiable surrogate of the twin} (L1) unlocks nonlinear fleet generation (L2),
the two physics-injection mechanisms H4 left untested --- exact-gradient physics losses and
twin-residual features --- (L6), and fast analysis-by-synthesis diagnosis. Mission-mix
generation (L2) then feeds the report-schedule findings of \cref{ch:tomography} with data
rather than geometry alone. This chapter reports the ten lines executed to date: L1 (the twin), L2 (the nonlinear
fleet), L6 (H4 re-opened), L5 (architecture breadth, \cref{sec:f8-l5}), L4 (the recoverable
fraction, \cref{sec:f8-l4}), L9 (the PCS metric, \cref{sec:f8-l9}), L7 (sensor drift,
\cref{sec:f8-l7}), L-H2 (real vs virtual sensors, \cref{sec:f8-lh2}), L-RUL (an advanced
classical prognostic, \cref{sec:f8-lrul}) and L-ICM, which turns the study's own premise ---
that the ICM's \emph{structure} outlives the twin's calibration error --- into a measurement
(\cref{sec:f-icm}). The heavier remaining lines --- real-station
EGT, N-CMAPSS DS02 and map-shape calibration --- queue behind those, each a pre-registrable
study in its own right.
Every line inherits the house discipline: a hypothesis frozen before the run, and a gate that
is allowed to fail.

\section{L1: a differentiable neural twin --- gate passed}
\label{sec:f8-l1}

\paragraph{Why.} The fleet generator is a linearization --- audited to be at the noise
floor for fleet-realistic health states (\cref{sec:audit-nonlin}), but a declared
approximation nonetheless, and one whose error grows exactly where interesting things
happen (end of life, large faults). The full pyCycle model is exact but costs seconds per
snapshot and offers no gradients. The goal: a model with pyCycle's fidelity at
microseconds per evaluation, differentiable end to end.

\paragraph{How.} 2\,400 full pyCycle solves sampled across the design range (all ten
health parameters in $\pm3.5\,\%$, commanded N1 across the cruise band; 8 parallel
workers, zero failed solves, $\sim$14 minutes on the M5). The architecture that passed is
the house pattern at model scale: \textbf{the linearization plus a learned residual} --- a
three-layer multilayer perceptron (MLP)---a plain feed-forward neural network---corrects
what the linear model misses, inheriting its smoothness where
physics is nearly linear. A pure MLP was tried first and could not match the
linearization's precision on N2; the residual form fixed that immediately.

\paragraph{The gate, including its two wrong drafts.} The acceptance gate went through
three disclosed revisions --- each a lesson in evaluation design. Draft one compared the
surrogate's \emph{maximum} error against the linearization's \emph{P95} (a statistic
mismatch). Draft two fixed the statistics but kept the fleet audit's numbers as the
baseline --- the \emph{wrong population}: the audit measured small, fleet-realistic health
states, while the surrogate's test set spans the full $\pm3.5\,\%$ design range, where the
linearization is an order of magnitude worse (N2 nonlinear residual: 0.42\,\% there versus
0.035\,\% in the audit). The final gate evaluates both models \emph{on the same held-out
test set} and demands strict like-for-like improvement.

\begin{figure}[htbp]
  \centering
  \includegraphics[width=0.62\textwidth]{l1_surrogate}
  \caption{L1 gate: median error against full pyCycle per cockpit channel, on the same
  held-out solves --- the surrogate (blue) is an order of magnitude closer than the
  linearization (gray). Auto-generated.}
  \label{fig:surrogate}
\end{figure}

\paragraph{Result.} \Cref{fig:surrogate}: the surrogate beats the linearization by
13--19$\times$ on every cockpit channel --- N2 at 0.008\,\% median where the linear model
sits at 0.12\,\%, EGT at 0.03\,\% versus 0.54\,\%. In exchange for 14 minutes of solve
generation, the project now owns a twin that evaluates in microseconds, differentiates
exactly, and is more faithful to the physics than the generator that produced every result
so far --- which is precisely what makes the next lines possible.

\section{L2: the nonlinear fleet (SynCFM56 v2)}
\label{sec:f8-l2}

With both family surrogates gated (\cref{sec:f8-l1}), the fleet generator's physics
emitter becomes swappable: setting \code{fleet.emitter = surrogate} replaces the
linearization of \cref{eq:generator} with the neural twin, at the same
40-second-per-fleet cost (CPU inference inside the existing per-engine workers). Everything
else --- trajectories, sensor model, seeds, splits --- is byte-identical, so v2 differs
from v1.1 in exactly one controlled variable: the fidelity of the true channels.
\textbf{v1.1 is left frozen} (its hashes anchor the F5/F7 pre-registrations); v2 lives
alongside it as \code{data/processed/fleet\_v2/}.

Two re-audits confirm v2. First, fidelity: re-solving full pyCycle at 60 states sampled from
the v2 fleet and comparing both generators against it (\cref{fig:v2-fidelity}), the
surrogate is 3--7$\times$ closer to true physics than the linearization on every cockpit
channel --- EGT median error $0.04\,\%$ versus $0.31\,\%$. Second, the difficulty gate:
the trivial classifier scores $58.0\,\%$ on v2's acute-isolation task (versus $58.1\,\%$
on v1.1), safely below the $60\,\%$ ceiling --- the nonlinear physics did not accidentally
make the benchmark easier, the one risk that mattered. Realism is intact (wash sawtooth
\SI{+2.9}{\degreeCelsius} recovery, end-of-life EGT and fuel-flow signs correct in $100\,\%$ of
engines, measured noise matching the catalog).

\begin{figure}[htbp]
  \centering
  \includegraphics[width=0.62\textwidth]{l2_fidelity}
  \caption{L2 fidelity: median generator error against full pyCycle, per cockpit channel ---
  the v2 surrogate (blue) is several times closer than the v1 linearization (gray).
  Auto-generated.}
  \label{fig:v2-fidelity}
\end{figure}

\section{L6: H4 re-opened --- and refuted a second way}
\label{sec:f8-l6}

The F5 stacking hybrid failed by concatenating smeared GPA \emph{estimates}
(\cref{sec:res-h4}). L6 tries the mechanism that failure suggested: not the estimate but
the \emph{residual}. Per snapshot, the linear WLS residual
$\mathbf{r} = (\mathbf{I} - \mathbf{H}_{\text{ref}}\mathbf{M})\,\Delta\mathbf{z}$
(with $\mathbf{M}$ the regularized WLS operator, \cref{eq:wls}) is the part of the
measurement the linear GPA \emph{cannot} explain. On v1 this was near zero by construction
(generator and GPA shared the linearization); on v2 the generator is the nonlinear
surrogate, so the residual carries real information --- the twin-residual feature. The RUL
GRU input grows from four channels to seven, and the experiment is frozen exactly as F5's
H4 was (\code{docs/prereg-v3.md}, tag \code{prereg-v3}): pure vs hybrid at 10/25/100\,\%
of training engines, three seeds, single confirmatory pass, thresholds fixed \emph{before}
running.

\begin{figure}[htbp]
  \centering
  \includegraphics[width=0.62\textwidth]{l6_hybrid}
  \caption{H4 on v2: the twin-residual hybrid (red) is worse than the pure GRU (blue) exactly
  where it should help most --- scarce data --- and no better at full data. Auto-generated.}
  \label{fig:h4v2}
\end{figure}

\Cref{fig:h4v2}: the hybrid is \emph{worse} where it was supposed to help most (2\,099 vs
1\,944 cycles at 10\,\% data, 1\,268 vs 1\,233 at 25\,\%) and statistically
indistinguishable at full data; the Wilcoxon test runs the wrong way ($p = 0.99$ against the
hybrid). \textbf{H4 is now refuted a second, independent way.} Two different physics-injection
mechanisms --- estimate-stacking and residual-augmentation --- both fail to improve RUL, on
both the linear and the nonlinear fleet. The transferable finding hardens: for this
prognosis task the raw corrected deviations already contain the remaining-life signal, and
bolting on physics-derived features adds variance, not information. Under the frozen
protocol, mechanism M3 (a physics-constrained loss) is therefore \emph{deferred}: with two
negatives in hand, the burden of proof has shifted, and any future M3 attempt should itself
be pre-registered against this record.

\paragraph{The differentiable twin's real payoff.} L6's negative does not diminish L1: the
surrogate's value is as a fast, exact-gradient \emph{generator} and analysis-by-synthesis
tool (L2; the sensor-drift line is closed in \cref{sec:f8-l7}), not as a feature
factory for a learner that already had the signal.

\section{L5: is the RUL advantage architecture-robust?}
\label{sec:f8-l5}

The headline prognosis result (H3) used a GRU. Is the advantage a property of \emph{learning
from sequences}, or an accident of that one architecture? L5 answers it under the fairness
rule that governs the whole project: three sequence architectures --- the GRU, a temporal
convolutional network (dilated causal convolutions) and a Transformer (self-attention) ---
get \emph{identical} treatment (same inputs, data, RUL cap, epoch budget, three seeds, and
fixed sensible hyperparameters). This is an equal-modest-effort comparison, disclosed as
such (\code{docs/prereg-v8.md}, tag \code{prereg-v8}); it is deliberately \emph{not} a
per-architecture tuning campaign, which would reintroduce the effort asymmetry the project
exists to avoid.

\textbf{H5L.1 confirmed}: all three architectures beat the tuned traditional baseline at
$90\,\%$ life (\cref{fig:arch}), and by nearly the same margin --- RMSE $827$--$888$ cycles
against the traditional $1\,981$, a $\sim$2.3$\times$ improvement whichever network is used.
The convolutional and attention models land within a few percent of the recurrent one; the
best (TCN) beats it slightly, but the finding is precisely that architecture choice barely
matters. The RUL advantage is a property of learning from the deviation sequence, not of the
GRU --- which hardens H3 rather than qualifying it.

\begin{figure}[htbp]
  \centering
  \includegraphics[width=0.72\textwidth]{l5_architectures}
  \caption{L5: RUL RMSE (mean of three seeds) for three AI architectures under equal effort;
  all sit far below the tuned traditional baseline at 90\,\% life (dashed). The advantage is
  the learning, not the architecture. Auto-generated.}
  \label{fig:arch}
\end{figure}

\section{L4: how much will a wash recover?}
\label{sec:f8-l4}

A fleet manager washing a compressor wants to know the payoff \emph{before} washing: what
share of an engine's lost EGT margin is fouling (which a wash restores) versus permanent wear
(which it does not). The composable ground truth (\cref{sec:degradation-models}) makes this a
supervised target: replaying each engine's per-mechanism contributions and projecting them
through the takeoff EGT row gives the true recoverable fraction at every cycle. The question
(\code{docs/prereg-v7.md}, tag \code{prereg-v7}): can it be read off the measured deviation
trajectory alone?

\textbf{H4L.1 confirmed, decisively.} A histogram gradient-boosting regressor on
trailing-window shape features (level, slope, curvature and sawtooth roughness of the cockpit
deviations) predicts the recoverable fraction with test $R^2 = 0.86$ and Spearman
$\rho = 0.86$ ($p < 10^{-4}$), against a mean-predictor baseline of zero (\cref{fig:recoverable}).
The mechanism is honest and physical: fouling betrays itself through the wash sawtooth its
recoveries leave in the deviation history, so an engine's own past wash responses reveal how
much the next wash will buy. This is the batch's first clean positive limitation-overcoming
result --- and an actionable one, converting the composability of the synthetic ground truth
into a wash-planning tool an operator could use.

\begin{figure}[htbp]
  \centering
  \includegraphics[width=0.44\textwidth]{l4_recoverable}
  \caption{L4: predicted versus true recoverable fraction of lost EGT margin on the test
  split (dashed = perfect). The estimator reads the wash-restorable share off the deviation
  trajectory with $R^2 = 0.86$. Auto-generated.}
  \label{fig:recoverable}
\end{figure}

\section{L9: can the Physics-Consistency Score be validated?}
\label{sec:f8-l9}

Contribution C5, the Physics-Consistency Score (\cref{sec:pcs}), was a null result on the
weak v0 classifier. L9 asks the fair question under pre-registration
(\code{docs/prereg-v5.md}, tag \code{prereg-v5}): when a classifier reasons more physically,
does PCS rise? Three diagnosis classifiers are scored --- one on the \emph{extended} sensor
set (where the ICM makes faults more separable, \cref{sec:bt-honesty}), one on the cockpit
set (the v0 setting), and a negative control trained on \emph{shuffled} labels
(\cref{tab:pcs}).

\begin{table}[htbp]
  \centering\small
  \caption{L9 PCS validation: isolation accuracy and mean PCS for a competent, a confused,
  and a label-shuffled classifier. Auto-generated.}
  \label{tab:pcs}
  % AUTOGENERATED by scripts/make_report_assets.py — do NOT edit by hand
\begin{tabular}{lrr}
\toprule
Classifier & Isolation acc. & Mean PCS \\
\midrule
  competent (extended sensors) & 0.39 & 0.251 \\
  confused (cockpit) & 0.35 & 0.180 \\
  control (shuffled labels) & 0.00 & 0.205 \\
\bottomrule
\end{tabular}

\end{table}

The verdict is split and, characteristically, deflating in an informative way.
\textbf{H9.2 confirmed}: for the extended classifier PCS is higher on the episodes it gets
right ($0.27$) than on those it gets wrong ($0.24$), so the metric does carry a little
signal about physical correctness. But \textbf{H9.1 is refuted as frozen}: the competent
classifier's PCS ($0.25$) does not significantly exceed the shuffled control's ($0.21$,
Mann--Whitney $p = 0.17$), and the control even edges the honest cockpit classifier. The
reason is the deeper finding, not a bug: \emph{no available classifier is strongly physical},
because acute isolation stays hard even on the extended set (accuracy $0.39$) --- the H2 wall
again. Without a genuinely physical reasoner to validate against, PCS cannot be shown to be
more than weakly informative here. C5 is therefore only partially validated; a clean test
would need a positive control (e.g.\ the extended-sensor WLS estimate, physical by
construction), which is left as pre-registrable future work rather than claimed post hoc.

\section{L7: the lying thermocouple, and the limit of a smarter filter}
\label{sec:f8-l7}

Case C (\cref{ch:cases}) showed a drifting EGT sensor corrupting every downstream method. The
classical remedy is an \emph{augmented-state} Kalman that estimates the sensor bias $b$
alongside the health vector --- the $\mathbf{S}\mathbf{b}$ term of \cref{eq:gpa-linear} made
into a tracked state (\code{docs/prereg-v6.md}, tag \code{prereg-v6}). Two questions: does it
recover the drift, and does it stop the drift from poisoning the diagnosis?

\textbf{It recovers the drift.} The augmented estimate of $b$ tracks the true bias with a
median Spearman $\rho = 0.83$ ($p = 0.004$) across the thirteen drifting-EGT engines. (The
\emph{frozen} confirmatory was restricted to the test split, which contained only two such
engines --- a pre-registration design limitation, disclosed --- so the significant number is
the disclosed exploratory one over all splits; the per-engine recovery is strong either way.)

\textbf{But it barely cleans the diagnosis.} The spurious health error the drift induces on
the hot-section parameters falls only marginally, from $1.69\,\%$ to $1.64\,\%$ ($\sim3\,\%$).
On the cockpit rank-3 measurement set the bias stays confounded with the health directions it
mimics, so the filter can \emph{track} the drift but cannot \emph{un-corrupt} the estimate.
The lesson is the one Case C already implied and the certificate (\cref{ch:breakthrough})
formalizes: against a lying sensor, a smarter algorithm buys awareness, not accuracy ---
the fix is instrumentation (cross-channel redundancy or the extended set), not filtering.

\section{Breaking the wall --- and why only a real sensor can}
\label{sec:f8-lh2}

The H2 refutation (\cref{sec:res-h2}) left an actionable question the certificate answered
geometrically (\cref{sec:bt-honesty}): the confusable wall is informational, so real
gas-path sensors should break it. This line tests that directly, and pairs it with a sharper
question --- can a \emph{virtual} sensor, one the twin predicts from the cockpit data an
operator already has, do the same for free? Under pre-registration (\code{docs/prereg-v9.md},
tag \code{prereg-v9}) the traditional nearest-signature isolator runs on the confusable test
episodes under three sensor conditions: cockpit only; cockpit plus \emph{model-predicted}
station channels; and cockpit plus the \emph{real} station channels (\cref{fig:wall}).

The result is as clean as the theory (\cref{eq:gpa-linear}) demands. Real station sensors take
confusable isolation from $0.15$ to $0.92$ --- $+77$ percentage points, McNemar $p = 0.001$:
\textbf{H-H2.1 confirmed}, the wall breaks the moment the missing information is physically
measured. The virtual sensors leave it at exactly $0.15$ ($p = 1.000$): \textbf{H-H2.2
confirmed}, they add nothing. This is the data-processing inequality made concrete --- a
quantity computed from the cockpit data cannot carry information the cockpit data did not
already hold, so it cannot separate faults the cockpit cannot. Only a probe placed where the
two faults differ --- the HPC exit --- supplies the missing dimension.

\begin{figure}[htbp]
  \centering
  \includegraphics[width=0.5\textwidth]{wall_break}
  \caption{Confusable-fault isolation accuracy under three sensor conditions. A virtual
  (twin-predicted) station sensor matches the cockpit exactly; only a real station sensor
  breaks the wall. Auto-generated.}
  \label{fig:wall}
\end{figure}

The industrial reading is precise and, in this case, unusually cheap: the station parameters
that break the wall (PS3, T25, T3) are already measured by the CFM56-7B FADEC for engine
control, so the barrier to using them for diagnosis is recording and downlinking that data,
not installing new hardware. And the negative half is a purchasing safeguard: no software vendor,
however sophisticated, can substitute a virtual sensor for the missing measurement --- the
certificate of \cref{ch:breakthrough} says exactly what the current sensors can and cannot
know, and this line shows that boundary is real.

\section{L-RUL: does an advanced classical prognostic narrow H3?}
\label{sec:f8-lrul}

The traditional RUL of \cref{ch:results} used linear Theil--Sen extrapolation --- operational
practice, not the research frontier (\cref{sec:math-theilsen}). So part of the H3 gap could be
``AI beats a linear extrapolator'' rather than ``AI beats the best classical prognostic.''
This line settles how much, under pre-registration (\code{docs/prereg-v10.md}, tag
\code{prereg-v10}), by pitting Theil--Sen against a \emph{similarity-based} prognostic --- the
C-MAPSS-standard classical method, which matches a test engine's recent degradation curve
against the run-to-failure library of the training engines and so can follow the nonlinear,
accelerating shape a line cannot (\cref{fig:lrul}).

Both hypotheses confirm, and the honest reading is a refinement, not a reversal.
\textbf{H-RUL.1 confirmed}: the advanced method nearly halves the traditional error at
$90\,\%$ life, $1\,981 \to 1\,118$ cycles --- much of the linear extrapolator's deficit was
its inability to bend. \textbf{H-RUL.2 also confirmed}: the AI still wins, $858$ versus
$1\,118$ (Wilcoxon $p = 0.002$). So H3 survives against a genuinely strong classical baseline,
but its margin should be stated honestly at two scales: the AI is $\sim$2.3$\times$ better than
the RUL an operator fields \emph{today} (Theil--Sen), and $\sim$1.3$\times$ better than the
best \emph{advanced} classical prognostic. The learning advantage in prognosis is real ---
the GRU captures late-life curvature and cross-channel structure the similarity match does
not --- but it is a $30\,\%$ edge over a good classical method, not an order of magnitude.

\begin{figure}[htbp]
  \centering
  \includegraphics[width=0.52\textwidth]{lrul_compare}
  \caption{RUL error at 90\,\% life for the operational (Theil--Sen), advanced classical
  (similarity-based) and AI prognostics. The advanced method closes most of the gap to the
  AI, which still leads. Auto-generated.}
  \label{fig:lrul}
\end{figure}

\subsection{The fraction sweep the first pass owed}
\label{sec:f8-lrul-sweep}

\paragraph{Why this correction was needed.} The paragraph above adjudicated AI against the
advanced classical at \emph{one} life fraction, $90\,\%$, while the headline H3 factor
($2.3$--$4.4\times$) is quoted against the operational Theil--Sen at \emph{three}. Comparing a
strong baseline at one point and a weak one everywhere is an asymmetry that flatters the AI,
so this pass (\code{prereg-v15}) runs the paired comparison at every evaluated fraction, one
Wilcoxon each, Holm-corrected across the three. The result changes the story materially, and
in the direction that costs the project a claim.

\begin{center}\small
\begin{tabular}{@{}lrrrrl@{}}
  \toprule
  Life fraction & Similarity & AI & AI margin & Holm $p$ & Verdict \\
  \midrule
  $50\,\%$ & $1\,998$ & $1\,689$ & $1.18\times$ & $0.41$ & no significant difference \\
  $70\,\%$ & $1\,064$ & $1\,059$ & $1.01\times$ & $0.41$ & statistical tie \\
  $90\,\%$ & $1\,118$ & $834$  & $1.34\times$ & $0.006$ & \textbf{AI wins} \\
  \bottomrule
\end{tabular}
\end{center}

\keyidea{Against the best classical prognostic, the AI's advantage is not a level shift across
life --- it is confined to late life. At mid-life the two methods are indistinguishable
($1.01\times$, $p = 0.41$); only in the last tenth of life does learning pull clear
($1.34\times$, $p = 0.006$). The $2.3$--$4.4\times$ headline is a statement about the
\emph{operational} baseline, and only that.}

\textbf{H-RUL.3 refuted} as a uniform claim: the AI beats the advanced classical at one of
three fractions, not three of three. Two readings follow, and both matter. Scientifically,
the learning advantage in prognosis is a \emph{late-life} phenomenon --- consistent with the
mechanism already identified, that the GRU captures the accelerating curvature and
cross-channel structure that appear once degradation dominates, and with F11's finding that
mid-life RUL spread is $87\,\%$ irreducible anyway (\cref{sec:bt-prognostic-floor}): where
there is nothing to win, nobody wins. Commercially, an operator whose classical prognostic is
already similarity-based should expect the AI to buy them the shop-visit horizon, not the
mid-life trend --- and should discount any vendor factor quoted against linear extrapolation.

\begin{figure}[htbp]
  \centering
  \includegraphics[width=0.62\textwidth]{lrul_sweep}
  \caption{AI against the \emph{advanced} classical prognostic at all three evaluated life
  fractions, with the Holm-corrected paired $p$ above each pair. The two methods are
  indistinguishable until the last tenth of life. Auto-generated (\protect\code{prereg-v15}).}
  \label{fig:lrul-sweep}
\end{figure}

\paragraph{Run-to-run disclosure.} The AI figures here ($1\,689/1\,059/834$) come from a
re-execution of the frozen configuration and differ from the confirmatory pass
($1\,694/1\,042/858$) by $0.3$--$3\,\%$; GRU training on the Metal backend is not
bit-reproducible across runs. The re-run numbers are used throughout this subsection so that
the pairing is internally consistent; the frozen H3 verdict is untouched. No conclusion here
turns on a difference that small --- the $70\,\%$ tie is a tie under either set.

\section{L-UQ: how much of the H5 interval advantage is the model, and how much is calibration?}
\label{sec:f-uq}

\paragraph{The confound.} H5 reported the AI's conformal $90\,\%$ interval at
$\pm1\,957$ cycles against the traditional slope band's $\pm11\,509$ --- a
$5.9\times$ narrower maintenance window (\cref{sec:res-h3h5}). But the two intervals were not
built the same way. The AI's was \emph{calibrated} by split conformal to hit $90\,\%$; the
traditional one was the raw Theil--Sen slope-percentile band, calibrated to nothing, and it
empirically covered $98.3\,\%$. An interval that over-covers is wide by construction. So the
$5.9\times$ mixes a genuine question --- is the AI's point prediction more accurate? --- with
one that has nothing to do with AI: was the interval calibrated at all. Conformal prediction
is model-agnostic; this project proved so itself by wrapping it around the classical
Kalman/CRB ellipsoid to repair H10.2 (\cref{sec:bt-honesty}). Whatever conformal contributes,
the traditional side can take for free.

\paragraph{The fix.} Wrap the identical split-conformal procedure around all three point
predictors --- AI GRU, similarity-based classical, Theil--Sen --- calibrating each on the
validation engines and evaluating each on the test engines
(\code{scripts/f\_uq\_reattribution.py}, \code{prereg-v15}). Every method then stands at the
same nominal $90\,\%$, and the width that remains is the width the \emph{predictor} earns.

\begin{center}\small
\begin{tabular}{@{}lrrrr@{}}
  \toprule
  Method (all conformalized at $90\,\%$) & Coverage & Half-width & vs AI & RMSE \\
  \midrule
  AI GRU                             & $0.883$ & $2\,028$ & $1.00\times$ & $1\,261$ \\
  Similarity (advanced classical)    & $0.917$ & $2\,717$ & $1.34\times$ & $1\,458$ \\
  Theil--Sen (operational)           & $0.767$ & $4\,627$ & $2.28\times$ & $3\,612$ \\
  \midrule
  \emph{frozen H5 traditional band (uncalibrated)} & $0.983$ & $11\,509$ & $5.88\times$ & --- \\
  \bottomrule
\end{tabular}
\end{center}

\keyidea{Held at equal coverage, the AI's maintenance window is $1.34\times$ narrower than the
best classical method's and $2.3\times$ narrower than fielded practice --- not $5.9\times$.
Roughly three quarters of the headline interval advantage was conformal calibration, which any
method can adopt, and one quarter was the model. The $1.34\times$ is the honest AI figure, and
it agrees with the late-life accuracy margin of \cref{sec:f8-lrul-sweep}.}

\begin{figure}[htbp]
  \centering
  \includegraphics[width=0.62\textwidth]{uq_reattribution}
  \caption{$90\,\%$ RUL interval half-widths with the empirical coverage each achieves. The
  three left bars give every predictor the same split-conformal calibration; the right bar is
  the frozen H5 traditional band, which is wide because it over-covers, not because the
  predictor is that much worse. Auto-generated (\protect\code{prereg-v15}).}
  \label{fig:uq-reattribution}
\end{figure}

\textbf{H-UQ.1 confirmed, narrowly}: the AI interval is still the tightest at honest coverage,
and the ordering of the three methods is unchanged. What does not survive is the magnitude.
Two caveats are owed. First, the Theil--Sen row under-covers ($0.767$ against nominal $0.90$):
its residuals are not exchangeable between validation and test --- the failure mode of a
predictor whose error distribution changes shape with life fraction --- so $2.28\times$ is a
\emph{lower bound} on the width it would need for honest coverage, and the operational
comparison is the more favourable one for the AI, not the less. Second, the AI's own coverage
sits at $0.883$, marginally below nominal; with $20$ calibration engines the realised coverage
of a split-conformal band still moves by a few points either way.

\paragraph{What an organization should take.} Ask a vendor for coverage before asking for
width. A narrow interval at unstated coverage is not a claim; a wide interval that over-covers
is not a failure. And if the incumbent classical prognostic is uncalibrated, the first and
cheapest upgrade is a conformal wrapper around \emph{it} --- perhaps $50$ lines of code, no
model change --- which captures most of the interval improvement before any learning is
purchased.

\section{L-MECH: the wall is instantaneous --- mechanism attribution over a history}
\label{sec:f13-mech}

\paragraph{Why the wall might not be the last word.} Everything this document has established
about the isolation limit follows from one structural fact: classical gas-path analysis is a
\emph{per-snapshot inverse solver}. The cockpit influence matrix has rank 3 against 10
unknowns, the confusable pair sits $1.3^{\circ}$ apart, halving every sensor $\sigma$ does not
separate it (\cref{sec:noise-sweep}), and the certificate stamps three efficiencies
unobservable even after accumulating Fisher information over an engine's whole flown history
(\cref{ch:breakthrough}). None of that is a statement about \emph{time}. It is a statement
about an \textbf{instantaneous estimand}.

The generator's chronic degradation is not instantaneous. It is five mechanisms with sharply
different temporal signatures (\cref{sec:degradation-models}): fouling saturating with a wash
sawtooth, erosion linear, clearance bilinear with a break-in knee, hot-section accelerating
with efficiency down \emph{and flow capacity up}, LPT wear linear. Per engine the true latent
is \textbf{five scalars}, not a free ten-vector per cycle --- a far smaller object than the one
the certificate declares unrecoverable.

\paragraph{The test.} Predict each mechanism's share of the lost EGT margin from the cockpit
deviation trajectory alone, at random cuts on engines held out by ID, with ground truth
replayed deterministically from each engine's seed. Two families under a \textbf{symmetric
25-trial Optuna budget}, selected on validation and reported on test
(\code{scripts/f13\_gate1\_mechanism.py}, \code{prereg-v16}): physics-designed shape
descriptors with gradient boosting, and a sequence model over the resampled history.

\begin{center}\small
\begin{tabular}{@{}lrrrl@{}}
  \toprule
  Mechanism & truth mean/sd & $R^2$ hand & $R^2$ sequence & temporal signature \\
  \midrule
  \code{hot\_section} & $0.47/0.13$ & $\mathbf{0.780}$ & $0.648$ & accelerating power law \\
  \code{clearance}    & $0.22/0.09$ & $0.286$ & $\mathbf{0.735}$ & bilinear break-in knee \\
  \code{fouling}      & $0.08/0.04$ & $0.451$ & $\mathbf{0.512}$ & wash sawtooth \\
  \code{lpt\_wear}    & $0.04/0.01$ & $0.217$ & $0.190$ & plain linear \\
  \code{erosion}      & $0.09/0.03$ & $0.130$ & $-0.018$ & plain linear \\
  \midrule
  \emph{mean}         &             & \emph{0.373} & \emph{0.414} & \\
  \bottomrule
\end{tabular}
\end{center}

\begin{figure}[htbp]
  \centering
  \includegraphics[width=0.72\textwidth]{f13_mechanism}
  \caption{Per-mechanism attribution from the deviation trajectory, on held-out engines, under
  a symmetric tuning budget. The dotted line is the $R^2 = 0.30$ recoverability mark.
  Auto-generated (\protect\code{prereg-v16}).}
  \label{fig:f13-mechanism}
\end{figure}

\keyidea{Three of five mechanisms are recoverable from trajectory shape --- including \code{hot\_section}, the mechanism that loads the confusable pair H2 could not break, at $R^2 = 0.78$. The identifiability wall is a property of the \emph{instantaneous} estimand, not of the history.}

\textbf{Gate 1a confirmed.} The question every operator actually asks at a shop visit is not
``what is the ten-dimensional health state now'' but ``which mechanism is consuming this
engine'' --- and that question has an answer where the first does not.

\textbf{Gate 1b refuted as specified}, and the way it fails is the more useful half. The
criterion was that the sequence model beat hand features on at least three of five mechanisms;
it manages two, at a mean gain of $+0.041$. This is neither ``hand features suffice'' nor
``learning wins''. It is a split that tracks \emph{how parameterizable each signature is}: the
learner wins by $0.449$ on the bilinear knee, which resists closed form, and loses by $0.132$
on the accelerating power law, which the hand set encodes directly as a quadratic term plus
cross-channel slope ratios. The two plain-linear mechanisms defeat both families.

\paragraph{What this result is not.} It is a statement about \emph{identifiability}, not about
maintenance planning. Asked which mechanism \emph{dominates} at the moment a workscope is
decided, this fleet answers \code{hot\_section} for 100 of 100 engines
(\cref{sec:limitations-conclusions}), so a shop-visit hit rate built on that question is scored
against a $100\,\%$ majority class and measures nothing. The recoverable quantity is
\emph{how much} of each mechanism, not \emph{which} --- and for the one share with a clear
operational reading, the washable fraction, \cref{sec:f8-l4} already reports $R^2 = 0.86$ with
a simpler estimator than anything here.

\paragraph{Two caveats, both against the learner.} The hand descriptors were written by someone
who had read the fault catalogue --- one feature per catalogued mechanism --- so that baseline
is stronger than a practitioner without the generator source would build. And selection was
thin: validation is 245 cuts from ten engines, and the val/test gap is large (val $0.166$ hand
against $0.507$ sequence; test $0.373$ against $0.414$). The test column is quoted; the
validation ranking is not. Disclosed: the gate thresholds were written before any run, but the
symmetric budget was added after an untuned first pass whose architecture verdict was discarded
for exactly the asymmetry \cref{sec:f-uq} corrects elsewhere.

\section{L-INST: instrument or engine?}
\label{sec:f15-instrument}

\paragraph{The question behind every wasteful removal.} A line engineer's first question is not
``which module'' but \emph{``is this an engine problem at all, or is my instrument lying?''} A
removal triggered by a drifting thermocouple is a wasteful removal, the false-alarm analog of
\cref{sec:ops-conversion}. The fleet already carries the experiment: 22 of the 100 engines have
a slow additive bias ramp on one sensor, with the channel and onset in ground truth.

At a snapshot the two are \textbf{perfectly confusable, provably}. The cockpit influence matrix
maps $\mathbb{R}^{10} \to \mathbb{R}^{3}$ \emph{surjectively}, so every deviation vector ---
including a pure single-channel bias --- has a health-state preimage. This is the same rank
argument that makes H2 unbreakable. \Cref{sec:f13-mech} having just shown that a history can
break an instantaneous confusion, the obvious question is whether it does so here.

\paragraph{Design, and what the data forces.} Only cockpit-\emph{visible} drift can count: the
fleet drifts four channels but the cockpit set is $\{N_2, W_f, \mathrm{EGT}\}$, so T25 and PS3
biases carry no information in the input at all. The 14 EGT and $W_f$ engines are the
positives; the 8 invisible-drift engines are held out as a \textbf{fraud check}. The frozen
splits put only three drifted engines in test --- which is precisely why L7's confirmatory pass
was under-powered and said so (\cref{sec:f8-l7}) --- so this line uses grouped five-fold
cross-validation over engine IDs with pooled out-of-fold scores, a disclosed deviation from the
frozen-split convention. Classical family: the largest augmented-Kalman bias magnitude swept
across candidate channels, reusing L7's estimator. Learned family: a bidirectional GRU over the
trajectory, three seeds per fold.

\begin{center}\small
\begin{tabular}{@{}lrrr@{}}
  \toprule
  Family & AUC & 95\,\% CI & fraud check \\
  \midrule
  augmented Kalman (classical) & $0.614$ & $[0.430,\ 0.789]$ & $0.433$ \\
  Bi-GRU (sequence)            & $0.463$ & $[0.286,\ 0.657]$ & $0.405$ \\
  \bottomrule
\end{tabular}
\end{center}

\begin{figure}[htbp]
  \centering
  \includegraphics[width=0.70\textwidth]{f15_instrument}
  \caption{Separating a drifting sensor from a real fault, 14 visible-drift engines against 78
  clean, bootstrap intervals grouped by engine. Crosses mark the fraud check: engines whose
  drift is on a channel the cockpit set cannot see. Auto-generated (\protect\code{prereg-v18}).}
  \label{fig:f15-instrument}
\end{figure}

\textbf{H15.8 refuted.} Neither family separates the two. Both intervals contain $0.5$, and
the learned family sits \emph{below} it --- worse than a coin.

\paragraph{Re-run with the pooling corrected, and the conclusion held.} The learned arm
originally used the same defective bidirectional read diagnosed in \cref{sec:f18-bidir}
(\code{o[:, -1]}, which for a backward pass has seen one sample). Repeating this study with
the correct read moved the classical arm not at all --- $0.614$ either way, as it must, since an
augmented-state Kalman shares none of that code --- and moved the learned arm from $0.524$
\emph{down} to $0.463$. So the defect had been mildly flattering the learner, and the null is
not an implementation artifact: with a correctly pooled model the separation is worse, not
better. The superseded figures are retained in the verdict JSON.

\paragraph{The fraud check is what makes the null credible.} The eight engines drifting on a
channel the cockpit cannot observe carry no signal by construction, and both families score
them like clean engines. A pipeline leaking labels would have scored those high too. The null
is a null, not a bug --- and this is the same self-check the certificate enables throughout
this document: physics says where information cannot be, and a method claiming to find it there
is caught.

\keyidea{L7 showed the augmented Kalman \emph{tracks} an EGT bias at $\rho = 0.83$. This shows tracking is not attributing: because the ICM is surjective, a genuine degradation also produces a nonzero apparent bias, so bias magnitude carries no evidence about which of the two is happening. The surjectivity argument, previously a statement about one snapshot, holds empirically over whole histories.}

\paragraph{The rule these two lines sharpen into.} \Cref{sec:f13-mech} and this section are the
same experiment on different fault families, and together they give one statement:
\textbf{trajectory shape separates what has shape}. The sequence model won on the break-in knee
and the wash sawtooth, lost on the two plain-linear mechanisms, and here loses again on a slow
additive ramp with no discriminating event. Four instances, two families, two fault
populations. A history breaks an instantaneous confusion only when something happens in it.

\paragraph{Power, stated.} Fourteen positives. The classical interval does not exclude a
moderate effect, so the honest claim is \emph{no detectable signal at this sample size}, not
proof of none. The point estimates are near chance and the learned family has none at all. The
uncomfortable consequence stands either way: this project's diagnostic output, like any
cockpit-only gas-path diagnosis, is conditional on instrument health it cannot itself verify.

\section{L-EXT: an external replication, and what single-run rankings hide}
\label{sec:f14-ext}

\paragraph{Why replicate someone else's paper.} Every correction in this document so far has
been this project auditing itself. An external check asks whether the pattern generalizes:
take a recent, peer-reviewed, well-written deep-learning PHM result and put it through the same
standards. The subject is \textcite{mukherjee2026bilstm}, which reports a bidirectional LSTM on
C-MAPSS FD001 at RMSE $14.121$ and $R^2 = 0.885$, best of fifteen models, against a weakest
comparator --- linear regression --- at $22.914$.

The protocol is reproduced exactly (\code{scripts/f14\_bilstm\_replication.py},
\code{prereg-v17}): 12 sensors selected by $|r| \ge 0.5$ with the cycle column dropped as
leakage, min-max scaling, a 30-cycle window, piecewise-linear RUL capped at 125, two stacked
Bi-LSTM layers, Adam at $10^{-3}$, 25 epochs, ModelCheckpoint and
\code{ReduceLROnPlateau(factor=0.01)}. Two deviations are disclosed because PyTorch cannot
express them: Keras applies recurrent dropout inside the recurrence, and Table 3's ReLU cell
activation is unavailable in the fused LSTM.

\begin{center}\small
\begin{tabular}{@{}lrrrrr@{}}
  \toprule
  Protocol & \multicolumn{5}{c}{FD001 test RMSE, five seeds} \\
  \midrule
  paper protocol (checkpoint on test) & $12.501$ & $84.550$ & $84.550$ & $84.550$ & $16.624$ \\
  clean validation split              & $84.550$ & $84.550$ & $13.328$ & $14.349$ & $18.853$ \\
  \bottomrule
\end{tabular}
\end{center}

\begin{figure}[htbp]
  \centering
  \includegraphics[width=0.70\textwidth]{f14_bilstm}
  \caption{Five seeds per protocol on FD001. Red points are runs whose output layer died and
  which predict zero for every unit. Auto-generated (\protect\code{prereg-v17}).}
  \label{fig:f14-bilstm}
\end{figure}

\textbf{$84.550$ is not a bad run.} It is exactly $\sqrt{\overline{y^2}}$ on the capped test
truth: those runs predict \emph{zero} for every engine. Five of ten collapse that way, and the
mechanism is specified by the paper itself --- Table 3 puts ReLU on the single-unit output
layer, so once it dies the gradient is zero, and Table 4's
\code{ReduceLROnPlateau(factor=0.01, min\_lr=10^{-5})} floors the learning rate before it can
recover. Recorded caveat: collapse probability depends on initialization, which differs between
Keras and PyTorch, so the \emph{mechanism} transfers but this \emph{rate} does not.

\textbf{H14.1 confirmed --- the result replicates.} Conditional on training at all, the paper
protocol gives $14.562 \pm 2.916$ against a reported $14.121$, and the best seed reaches
$12.501$. This is a reproducible paper.

\textbf{H14.2 confirmed, and mildly.} Moving the checkpoint from the reported test set to
twenty validation units carved out of train costs $14.562 \to 15.510$: a $0.95$-cycle,
$\sim$6\,\% penalty. Smaller than expected --- test-set model selection was not doing much
work here. (The pre-registered criterion compared pooled means, which under a bimodal outcome
mostly compares collapse counts; the trained-only pair is the substantive figure and the frozen
verdict is kept beside it.)

\textbf{H14.3 --- and the baselines invert.} Against this project's fielded health-index
projection ($39.284$ on FD001) the learner is $2.53\times$ better. But that projection scores
\emph{worse} than the paper's own linear regression at $22.914$: it is the operationally
realistic method, not the strongest classical one. The demanding figure is therefore
$\mathbf{1.48\times}$ against the paper's weakest comparator --- squarely in the $1.3$--$1.5$
band this document keeps arriving at once a baseline is competent
(\cref{sec:f8-lrul-sweep,sec:f-uq}).

\textbf{H14.4 refuted --- the clean finding.} The seed standard deviation among trained runs is
$2.94$ cycles. The paper ranks fifteen architectures on single runs with gaps of $0.120$ against
LSTM, $0.480$ against TCN and $0.603$ against GRU. \textbf{Every gap is inside a fifth of one
seed standard deviation.} The ranking is noise, and that conclusion is independent of the
collapse issue.

\keyidea{An external, peer-reviewed result reproduces --- and its headline conclusion, that one architecture is best of fifteen, does not survive running the same code five times. The margin against a competent baseline lands at $1.48\times$, the same place this project's own re-audits landed.}

\paragraph{Run-to-run disclosure.} As in \cref{sec:f8-lrul-sweep}, recurrent training on the
Metal backend is not bit-reproducible: a repeat of this study moved the best R1 seed from
$12.996$ to $12.501$ and the R1 trained-only mean from $14.810$ to $14.562$, while the R2
column was unchanged. The figures quoted are from the run that generated
\cref{fig:f14-bilstm}. Nothing here turns on differences of that size --- least of all H14.4,
whose conclusion needs only that the seed spread exceed $0.603$ cycles, which it does by a
factor of five.

\section{L-HYB: physics injection fails for prognosis and works for attribution}
\label{sec:f17-hybrid}

\paragraph{A third attempt, and why it was allowed one.} This document has refuted
physics-into-learner injection twice: H4 stacked GPA state \emph{estimates} as features and
lost (\cref{sec:res-h4}), L6 fed twin \emph{residuals} on the nonlinear fleet and lost again
(\cref{sec:f8-l6}). The stated diagnosis was that smeared estimates are noise-bearing,
mutually correlated features carrying nothing the learner did not already have, and that
``physics injection needs an information channel the raw data lacks.'' The F10 certificate is
such a channel: per engine and per direction it is computed from influence-matrix geometry and
the engine's \emph{actual flown conditions}, never from measured values, so a learner reading
the deviation trajectory cannot derive it. That distinction is what H15.3 tests.

The task is mechanism attribution (\cref{sec:f13-mech}), the one place with real signal;
re-running RUL would only repeat H4. Three arms share F13's selected architecture so the input
channel is the only variable, and the design was frozen as \code{prereg-v20} before the
confirmatory run --- ten seeds, a one-sided paired $t$-test on $C - A$, and \textbf{the same
test on $B - A$ as a control that must not fire}.

\begin{center}\small
\begin{tabular}{@{}llrrr@{}}
  \toprule
  Arm & Input & ch & seed-mean $R^2$ & ensembled \\
  \midrule
  A & 4 deviation channels + age & 5 & $+0.248 \pm 0.107$ & $+0.365$ \\
  B & \quad + raw 10-dim Kalman GPA estimate & 15 & $+0.387 \pm 0.056$ & $+0.499$ \\
  C & \quad + identifiable-gated estimate + per-direction CRB & 25 & $+0.440 \pm 0.060$ & $+0.540$ \\
  \bottomrule
\end{tabular}
\end{center}

\begin{figure}[htbp]
  \centering
  \includegraphics[width=0.62\textwidth]{hybrid_arms}
  \caption{Per-seed attribution score for the three arms; grey lines join the same seed across
  arms, which is the pairing the statistics use. Auto-generated
  (\protect\code{prereg-v20}).}
  \label{fig:hybrid-arms}
\end{figure}

\textbf{H15.3 refuted by its own control.} $C - A = +0.191$ ($t = 5.48$, $p = 0.0002$), but
$B - A = +0.139$ ($t = 5.01$, $p = 0.0004$). The control fired. Both physics arms beat the pure
baseline decisively, so the gain is \emph{having the state estimate at all} --- channel count
--- not the certificate. The control existed precisely so this could not be claimed, and it
did its job.

\keyidea{The refutation exposes a bigger result: arm B \emph{won}. H4 and L6 refuted physics injection on the \textbf{prognosis} task; on \textbf{attribution} the identical injection helps decisively. Physics injection is not a free win and not a dead end --- it is task-dependent, and the task decides.}

\paragraph{Why the three results reconcile rather than contradict.} For remaining life, the
health-state estimate adds nothing over the raw deviations, because life here is driven by EGT
margin, which is already an input channel --- the learner had that information. For attributing
\emph{which} mechanism is degrading, the estimated health direction is exactly the relevant
physical quantity, and it remains so even when smeared across the confusable subspace: a
smeared estimate still says the loss is hot-section-ish rather than fan-ish. H4's verdict
should therefore be read as scoped to prognosis, not as a general statement about
physics-informed learning, and \cref{tab:f8-summary} now records both.

\paragraph{One post-hoc, disclosed and not claimed.} Whether certificate gating adds anything
\emph{on top of} the raw estimate is $C - B = +0.053$ ($t = 2.24$, $p = 0.026$). Nominally
significant, outside the frozen design, and it would not survive Holm correction against the
two pre-registered tests. Per mechanism it concentrates where a working certificate should act
--- \code{hot\_section} $0.621 \to 0.865$ and \code{clearance} $0.623 \to 0.785$, the
directions the certificate has most to say about --- while \code{fouling} and \code{lpt\_wear}
move the other way. Suggestive; the honest status is a pre-registrable hypothesis for a later
pass, not a result.

\paragraph{Disclosures.} The exploratory three-seed pass (arm means $0.323/0.355/0.412$, paired
$C-A$ of $+0.088$ at $t = 2.32$) was \emph{inconclusive} against its own criterion, and it is
what set the confirmatory direction; it is retained in the verdict JSON so the confirmatory
numbers cannot be read without it. Hyperparameters were selected by F13 for arm A's five
channels and inherited by B and C, so arms with more channels were never separately tuned ---
a limitation that cuts against C, not for it.

\section{L-CERT: the certificate contrast, with channel count held fixed}
\label{sec:f19-cert}

\paragraph{Chasing this project's own loose end.} \Cref{sec:f17-hybrid} reported a post-hoc
$C - B$ of $+0.053$ ($t = 2.24$, $p = 0.026$) --- whether gating the GPA estimate with the
certificate adds anything on top of the raw estimate --- and declined to claim it. Revisiting
it, that post-hoc \textbf{carried the same defect that had just killed the primary
hypothesis}: arm $C$ had 25 channels and arm $B$ had 15, so $C - B$ was capacity-confounded one
level down from $C - A$. Reporting it as evidence for the certificate would have repeated
exactly the error the control arm existed to catch.

Arm $C$ also varied two things at once: it \emph{masked} the estimate to the
certified-identifiable subspace \emph{and} appended the per-direction CRB as static channels.
Those are separate claims --- one about the certificate's tags, one about its magnitudes. Four
arms separate them, and two pairs are matched on channel count by construction
(\code{scripts/f19\_certificate\_isolated.py}, \code{prereg-v22}, fifteen seeds, paired):

\begin{center}\small
\begin{tabular}{@{}llrr@{}}
  \toprule
  Arm & Input on top of the four deviation channels & ch & mean $\pm$ sd \\
  \midrule
  $B$   & raw 10-dim Kalman GPA estimate & 15 & $+0.385 \pm 0.073$ \\
  $C_1$ & the estimate \emph{masked} to the identifiable subspace & 15 & $+0.267 \pm 0.121$ \\
  $C_2$ & raw estimate $+$ per-direction CRB channels & 25 & $+0.399 \pm 0.081$ \\
  $C$   & masked estimate $+$ CRB channels & 25 & $+0.406 \pm 0.066$ \\
  \bottomrule
\end{tabular}
\end{center}

\begin{center}\small
\begin{tabular}{@{}llrrl@{}}
  \toprule
  Role & Contrast & $\Delta$ & paired $t$ & $p$ (one-sided) \\
  \midrule
  \textbf{primary} (count-matched) & $C_1 - B$ & $\mathbf{-0.119}$ & $-3.72$ & $0.999$ \\
  secondary (count-matched) & $C - C_2$ & $+0.007$ & $0.24$ & $0.407$ \\
  \emph{reported, not a test} & $C_2 - B$ & $+0.014$ & --- & count-confounded \\
  \bottomrule
\end{tabular}
\end{center}

\begin{figure}[htbp]
  \centering
  \includegraphics[width=0.66\textwidth]{cert_isolated}
  \caption{The certificate contrast at fixed channel count; grey lines join the same seed across
  arms. Masking the estimate with the certificate is the red arm. Auto-generated
  (\protect\code{prereg-v22}).}
  \label{fig:cert-isolated}
\end{figure}

\textbf{H15.11 refuted, and in the opposite direction.} Masking with the certificate does not
help; it \emph{hurts}, by $0.119$ at $t = -3.72$, with channel count equal by construction so
no capacity explanation remains. And the CRB as an input channel contributes nothing:
$C - C_2 = +0.007$ ($p = 0.41$), also count-matched.

\keyidea{The L-HYB post-hoc was the channel count. Tested with the counts equalised, its sign \emph{inverts}. That is the second post-hoc in this chapter to dissolve once its confound is removed --- and it is why the control arm, not the headline, is the part of an experiment worth designing carefully.}

\paragraph{Why masking hurts, which should have been foreseen.} The certificate tags a direction
\code{unobservable} when its CRB is \emph{large} --- that is, where the estimate is
\emph{imprecise}, not where it is empty. Zeroing those components discards information that is
noisy but not absent: a smeared estimate of $\eta_{\mathrm{HPC}}$ still says \emph{hot-section-ish
rather than fan-ish}, and \cref{sec:f17-hybrid} is precisely the finding that a learner exploits
that. Masking destroys it, and the seed spread widens accordingly ($\pm0.121$ against $B$'s
$\pm0.073$).

So the certificate is right about what it certifies and wrong as a filter. \textbf{A precision
bound is not an absence of information} --- and using one to discard data is a misreading of the
instrument this project built. That is a genuine correction to how \cref{ch:breakthrough}'s
certificate should be used downstream, and it was only available because the contrast was run
with the confound removed.

\section{L-BIDIR: is the published architecture's advantage direction or capacity?}
\label{sec:f18-bidir}

\paragraph{The one architectural claim with an argument behind it.}
\Cref{sec:f14-ext} showed the source paper's fifteen-way ranking sits inside seed noise, but
that is a statistical objection, not an engineering one. The paper's actual subject ---
\emph{bi}directional recurrence --- does have a physical argument, and it predicts where it
should and should not pay:

\begin{itemize}
  \item \textbf{Attribution} is decided at removal, with the whole history in hand. A backward
    pass is then legitimate smoothing rather than forecasting, and should be informative: a
    late wash sawtooth re-labels ambiguous early fouling; an end-of-life flow-capacity rise
    re-labels a mid-life efficiency drop as hot-section creep.
  \item \textbf{Prognosis} runs the backward pass over an input window that is entirely past.
    No leakage, but no future context inside the window to exploit either.
\end{itemize}

Pre-registered as \code{prereg-v21}: bidirectionality should help attribution and barely help
prognosis. \textbf{The measurement inverted both halves.}

\paragraph{Separating the two things a Bi-LSTM changes.} A Keras bidirectional layer of
nominal width $w$ carries \emph{twice} the parameters of a unidirectional layer of width $w$,
so any bi-versus-uni comparison at equal nominal width confounds \emph{direction} with
\emph{capacity}. Three arms separate them, sharing everything else --- data, sequences, ten
seeds, paired on seed (\code{scripts/f18\_bidirectionality.py}):

\begin{center}\small
\begin{tabular}{@{}llrrl@{}}
  \toprule
  Task & Arm & mean $\pm$ sd & $\Delta$ vs uni & paired $p$ \\
  \midrule
  \multirow{3}{*}{attribution ($R^2$)}
    & \code{uni} (128)          & $+0.181 \pm 0.114$ & --- & --- \\
    & \code{bi\_equal} (64/dir) & $+0.230 \pm 0.097$ & $+0.049$ & $0.189$ \\
    & \code{bi\_double} (128/dir) & $+0.167 \pm 0.103$ & $-0.013$ & $0.601$ \\
  \midrule
  \multirow{3}{*}{prognosis (cy RMSE)}
    & \code{uni} (128)          & $1\,490.9 \pm 120.8$ & --- & --- \\
    & \code{bi\_equal} (64/dir) & $1\,475.6 \pm 172.9$ & $-15.3$ & $0.419$ \\
    & \code{bi\_double} (128/dir) & $\mathbf{1\,396.7 \pm 85.1}$ & $\mathbf{-94.2}$ & $\mathbf{0.016}$ \\
  \bottomrule
\end{tabular}
\end{center}

\begin{figure}[htbp]
  \centering
  \includegraphics[width=0.82\textwidth]{bidir}
  \caption{Direction against capacity on both tasks; grey lines join the same seed across arms,
  which is the pairing the statistics use. Auto-generated (\protect\code{prereg-v21}).}
  \label{fig:bidir}
\end{figure}

\keyidea{At \emph{equal} parameter budget, bidirectionality buys nothing on either task ($p = 0.19$ and $p = 0.42$). At \emph{doubled} budget it buys $94$ cycles of RUL error ($p = 0.016$) and still nothing on attribution. On a GRU cell, then, what a bidirectional layer contributes here is \textbf{capacity rather than direction}. Whether that reading extends to the LSTM the source paper actually uses is tested below --- and it does not survive.}

\textbf{H15.2 refuted, in both halves.} Bidirectionality did not help attribution, and it did
help prognosis. The stated mechanism was therefore wrong, and the likely reason is visible in
the design: the trajectory is resampled to a fixed 256 steps, so the \emph{forward} pass
already sees the global shape of the whole history. The backward pass adds little because the
discriminating structure is not localized at the start of the record. What the extra capacity
buys, it buys on prognosis --- the task that has to compress a whole degradation curve into one
number.

\paragraph{A cross-run instability worth stating.} The \code{uni} attribution arm scores
$+0.181$ here, $+0.248$ as arm A of \cref{sec:f17-hybrid}, and $+0.414$ in
\cref{sec:f13-mech} --- three runs of a nominally identical configuration. With a per-seed sd
of $0.114$ that is consistent, but the consequence is a rule for reading this chapter:
\textbf{only paired comparisons within a single run are reliable}, and numbers must not be
subtracted across studies. Every verdict above is a within-run paired test for exactly that
reason.

\subsection{The same test on an LSTM cell, which is what the paper uses}
\label{sec:f18-lstm}

\paragraph{An overreach of this document's own, caught and corrected.} Everything above uses a
\textbf{GRU} cell, while the source paper --- and \cref{sec:f14-ext}'s replication of it ---
uses an \textbf{LSTM}. The parameter-doubling argument is cell-agnostic in principle, but ``in
principle'' is not a measurement, and generalising from an untested cell is the same species of
error this chapter exists to correct. Repeating all three arms with LSTM cells, everything else
identical:

\begin{center}\small
\begin{tabular}{@{}llrrl@{}}
  \toprule
  Task & Arm & GRU & LSTM & LSTM $\Delta$ vs uni \\
  \midrule
  \multirow{3}{*}{attribution ($R^2$)}
    & \code{uni}        & $+0.181$ & $-0.220 \pm 0.225$ & --- \\
    & \code{bi\_equal}  & $+0.230$ & $-0.055 \pm 0.204$ & $+0.164$ ($p = 0.099$) \\
    & \code{bi\_double} & $+0.167$ & $-0.312 \pm 0.177$ & $-0.092$ ($p = 0.86$) \\
  \midrule
  \multirow{3}{*}{prognosis (cy)}
    & \code{uni}        & $1\,490.9$ & $\mathbf{1\,251.3 \pm 69.4}$ & --- \\
    & \code{bi\_equal}  & $1\,475.6$ & $1\,417.7 \pm 116.4$ & $+166$ \emph{worse} ($p = 0.995$) \\
    & \code{bi\_double} & $\mathbf{1\,396.7}$ & $1\,380.5 \pm 121.3$ & $+129$ \emph{worse} ($p = 0.997$) \\
  \bottomrule
\end{tabular}
\end{center}

\textbf{The finding does not transfer.} On a GRU the doubled-parameter bidirectional arm
\emph{improved} prognosis by $94$ cycles; on an LSTM it \emph{worsens} it by $129$. Opposite
sign. The claim that the published architecture's edge is capacity rather than direction is
therefore \textbf{restricted to the GRU cell} and is not supported for the cell the paper uses.

\paragraph{And this comparison cannot settle the question either.} The hyperparameters ---
hidden $128$, two layers, $\eta = 0.0057$, $122$ epochs --- were selected by \cref{sec:f13-mech}'s
Optuna search \emph{for a GRU}. All three LSTM attribution arms return \emph{negative} $R^2$,
which is not ``LSTM is poor at attribution'' but a model that is not converging: a learning rate
suited to a three-gate cell need not suit a four-gate one. So this experiment cannot separate
\emph{cell} from \emph{tuning}, and it is not evidence that bidirectionality hurts an LSTM.

\keyidea{Two architectures compared at hyperparameters tuned for one of them measure the tuning, not the architecture --- which is precisely the objection this chapter raises against a fifteen-way published ranking. The honest status is that the capacity-not-direction reading holds for GRU, is untested for LSTM, and would need a symmetric tuning budget per cell to become a claim about architecture at all.}

\paragraph{Disclosed: the first run of this line was void.} It read \code{o[:, -1]} from the
bidirectional layer. That position is correct for a forward pass but not for a backward one,
which starts at the end and has therefore seen a single sample there --- so half the
representation reaching the head was close to the raw last observation, and the bidirectional
arms were crippled by construction. It reported attribution collapsing from $+0.241$ to
$-0.008$, which would have looked like a dramatic refutation and was an implementation error.
The corrected read is forward-at-last concatenated with backward-at-first. The void numbers are
retained in the verdict JSON. The same defect was found and fixed in
\cref{sec:f15-instrument}'s learned arm.

\section{Gate T: would transients open the confusable angle?}
\label{sec:gate-t}

\paragraph{The proposal, and the question it did not ask.} A natural objection to this
document's central negative is that the influence coefficient matrix is a \emph{steady-state}
Jacobian, while an engine's transient response is governed by rotor inertia, heat soak and
volume dynamics --- physics that depends on health parameters differently. Two faults
confusable in a steady signature need not be confusable in a step response. The argument is
sound in general. What it never establishes is whether it holds for \emph{this} pair:
$\eta_{\mathrm{HPC}}$ and $\eta_{\mathrm{HPT}}$ are both hot-section-adjacent, and their
dynamic signatures may be as close as their steady ones.

Building a transient deck is a project in itself, so this gate prices it first
(\code{scripts/f16\_gate\_t\_transient.py}, \code{prereg-v19}). A transient contributes two
separable things, and only one is computable from what this project owns.

\paragraph{Part 1 --- the quasi-steady content, and its upper bound.} A transient sweeps the
engine through a range of operating points, and each point has its own measured ICM. Sampling
many points is exactly the mechanism F7 exploited (\cref{ch:tomography}, $+79\,\%$ from adding
one report condition). So the generous limit of that route is: sample the \emph{entire}
measured envelope --- all six calibrated points --- at a fixed total sample budget, and see how
far the angle opens. Signatures are whitened by $\mathbf{R}^{-1/2}$ first, so ``separable''
means separable relative to sensor noise.

\begin{center}\small
\begin{tabular}{@{}lrr@{}}
  \toprule
  Operating point & cockpit & extended \\
  \midrule
  takeoff        & $0.75^{\circ}$ & $25.74^{\circ}$ \\
  takeoff\_hot   & $0.78^{\circ}$ & $28.07^{\circ}$ \\
  climb          & $0.90^{\circ}$ & $26.24^{\circ}$ \\
  cruise         & $1.55^{\circ}$ & $26.37^{\circ}$ \\
  cruise\_39k    & $0.83^{\circ}$ & $24.86^{\circ}$ \\
  cruise\_lowpwr & $0.79^{\circ}$ & $29.76^{\circ}$ \\
  \midrule
  \textbf{whole envelope stacked} & $\mathbf{1.38^{\circ}}$ & $\mathbf{26.71^{\circ}}$ \\
  \bottomrule
\end{tabular}
\end{center}

\begin{figure}[htbp]
  \centering
  \includegraphics[width=0.72\textwidth]{gate_t}
  \caption{Confusable-pair angle at each calibrated operating point and over the whole measured
  envelope, noise-whitened, log scale. Auto-generated (\protect\code{prereg-v19}).}
  \label{fig:gate-t}
\end{figure}

\textbf{Gate T refuted on the quasi-steady route, decisively.} Stacking every operating point
the twin has been calibrated at gives $1.38^{\circ}$ --- \emph{below} cruise alone at
$1.55^{\circ}$, because the extra points pull toward the common direction rather than away from
it. The pair points nearly the same way at every condition in the envelope; the
$1.3^{\circ}$ of \cref{sec:icm} was never an artifact of choosing cruise.

The Cram\'er--Rao bound does improve on the spread budget ($1.73 \to 1.12\,\%$ for
$\eta_{\mathrm{HPC}}$, $1.49 \to 0.92\,\%$ for $\eta_{\mathrm{HPT}}$) --- more conditions means
more information --- but the \emph{direction} is unchanged, which is what confusability is
about. More precision along a shared axis does not separate two faults on that axis.

\keyidea{Whatever a transient is worth here, none of it comes from operating-point diversity. Sweeping the entire calibrated envelope leaves the confusable pair at $1.4^{\circ}$. The quasi-steady route is closed.}

\paragraph{Part 2 --- pricing the part that is not computable.} Spool time constants and
heat-soak lags are genuinely new measurements, not a resampling of the steady map, and they
need a quasi-steady transient deck this project does not have. Rather than invent one, the gate
inverts the question: appended to the whole-envelope cockpit stack, how discriminating would a
time-constant measurement have to be to reach a target angle? In whitened units --- where the
entire steady signature has norm $8.1$ --- the differential sensitivity required is
$\mathbf{0.7}$ for $5^{\circ}$, $1.4$ for $10^{\circ}$ and $2.9$ for $20^{\circ}$.

So a dynamic channel carrying about $9\,\%$ of the steady envelope's discriminating power
would be enough to break the pair. That is not a large number, and a transient study is
therefore \emph{not} ruled out --- but two cautions belong with it. First, the requirement is on
the \emph{difference} between the two parameters' dynamic sensitivities, and both act the same
way on the high-pressure spool: degrading either $\eta_{\mathrm{HPC}}$ or $\eta_{\mathrm{HPT}}$
reduces net accelerating power and lengthens the spool time constant. The difference is
therefore much smaller than either sensitivity, which is precisely the trap the steady case
already fell into. Second, $0.7$ whitened units means a $0.7\sigma$ shift in the estimated time
constant per $1\,\%$ of health change on a single transient; if the per-transient sensitivity
is nearer $0.1$--$0.2\sigma$, as a first-order spool balance suggests, then roughly
$25$--$50$ comparable transients must be averaged to reach it.

\paragraph{What the extended column says, and it is the practical answer.} The pair sits at
$25$--$30^{\circ}$ at \emph{every} operating point once station probes are present, and the
$\eta_{\mathrm{HPC}}$ bound tightens to $0.023\,\%$. The wall is a property of the
\emph{cockpit} sensor set, not of steady-state analysis as such. An operator choosing between
funding a transient-diagnostics programme and fitting station instrumentation has the answer
already, and it agrees with \cref{sec:f8-lh2}: real probes lift confusable isolation
$0.15 \to 0.92$, and nothing in this gate suggests transients would be the cheaper route to the
same place.

\section{L-EXT2: the wall, measured on an engine model this project did not build}
\label{sec:f20-ncmapss}

\paragraph{The objection this answers.} Every geometric claim in this document descends from one
object: the influence coefficient matrix of \emph{our} pyCycle twin. The $1.3^{\circ}$
confusable pair, the rank-3 cockpit limit, the certificate of \cref{ch:breakthrough} and the
negative verdict on H2 all follow from it. \Cref{sec:f-icm} showed the verdict survives
perturbing our own calibration --- but that is still our model auditing itself. The open
objection has always been sharper: \emph{is the wall a fact about turbofans, or an artifact of
one deck?}

N-CMAPSS answers it, because it is the only external run-to-failure dataset that ships the
\textbf{true health parameters} alongside the measurements. It is a different engine, simulated
with a different code (NASA's C-MAPSS~\cite{AriasChao2021}), by different people, with 14
sensors instead of 3.

\paragraph{Getting a Jacobian out of a black box.} No influence matrix is published for
C-MAPSS, and the benchmark literature treats the model as closed. But an ICM \emph{is} a
Jacobian, and every N-CMAPSS sample pairs a health state $\boldsymbol{\theta}$ with
measurements $\mathbf{x}_s$ at a known operating condition $\mathbf{w}$. Regressing sensor
deviations on $\boldsymbol{\theta}$ within bins of $\mathbf{w}$ estimates the matrix from data
--- the ordinary way to obtain a Jacobian when the model cannot be perturbed
(\code{scripts/f20\_ncmapss\_icm.py}, \code{prereg-v23}).

\textbf{The identifiability gate comes first, and it had to.} Degradation is coordinated within
a unit, so the $\boldsymbol{\theta}$ columns correlate and a single subset need not separate
them. The seven failure modes are spread \emph{across} subsets, so pooling is what buys the
variation. Over nine of the ten repository files (699\,012 subsampled rows, 430 condition bins):
condition number $\mathbf{5.6}$, effective rank $\mathbf{10/10}$. All ten health parameters are
separately identifiable, and the fit proceeds.

\begin{center}\small
\begin{tabular}{@{}lrr@{}}
  \toprule
  Sensor set & this study (pyCycle, CFM56-7B) & N-CMAPSS (C-MAPSS) \\
  \midrule
  cockpit class (3 channels)      & $1.30^{\circ}$  & $\mathbf{0.66^{\circ}}$ \\
  full instrumentation            & $26.71^{\circ}$ & $\mathbf{20.92^{\circ}}$ \\
  \bottomrule
\end{tabular}
\end{center}

\begin{figure}[htbp]
  \centering
  \includegraphics[width=0.62\textwidth]{ncmapss_angle}
  \caption{The $\eta_{\mathrm{HPC}}$ against $\eta_{\mathrm{HPT}}$ angle on two independent
  engine models, log scale. Cockpit-class channels are $\{N_2, W_f, \mathrm{EGT}\}$ here and
  their N-CMAPSS analogues $\{N_c, W_f, T_{48}\}$. Auto-generated
  (\protect\code{prereg-v23}).}
  \label{fig:ncmapss-angle}
\end{figure}

\keyidea{On an engine model built by other people, in other software, for a different engine, the confusable pair sits at $0.66^{\circ}$ on cockpit-class instrumentation and opens to $20.9^{\circ}$ with full instrumentation. Both the wall and its cure reproduce. \textbf{The central negative of this document is not an artifact of its own simulator.}}

\paragraph{What may and may not be concluded.} The two numbers are not equal and should not be
quoted as if they were. Different engines, and the sensor sets correspond only approximately ---
$T_{48}$ is an HPT-outlet total temperature where our EGT is a station-4.5 proxy
(\cref{sec:limitations}). What transfers is the \emph{order of magnitude and the structure}: a
pair separated by under two degrees on cockpit-class instrumentation in both models, opening by
a factor of twenty to thirty once station probes are present in both. That is precisely the
class of claim \cref{ch:conclusions} says is transferable --- rankings, mechanisms and geometry,
not absolute values --- and it is now demonstrated rather than asserted.

Two further honest limits. The matrix here is \textbf{estimated}, not exact: it carries
regression error, which is another reason to read the comparison at the level of magnitude.
And one of the ten files (\code{DS08d}) unpacked truncated by 32 bytes and was skipped, recorded
in the verdict rather than passed over silently; the two other all-component subsets
(\code{DS08a}, \code{DS08c}) loaded, so no failure mode is missing from the pool.

\paragraph{Scope.} This study sits \emph{outside} the reproducible pipeline. It needs a
14.68\,GB download, and \code{make all} must never require it --- the same disclosure discipline
applied to the FD001 substitution in \cref{sec:sim-to-real}. The gate reasoning and the download
provenance are recorded in \code{docs/f20-ncmapss-gate.md}.

\section{L-CERT2: porting the certificate and the floor to external truth --- and the control they never had}
\label{sec:f21-port}

\Cref{sec:f20-ncmapss} established the geometry externally. The obvious next step is the one
the plan wanted all along: recompute the two results this project could only validate against
its own generator --- the certificate (\cref{ch:breakthrough}) and the prognostic floor
(\cref{sec:bt-prognostic-floor}) --- on NASA's truth. One of them ports. The other does not, and
in failing it exposed a defect in the original.

\subsection{The floor ports, and confirms}

F11 needs no influence matrix at all: condition on each unit's \emph{true} health state, take
its nearest neighbours in health space, and read the spread of their true remaining lives. Over
99 N-CMAPSS units (\code{scripts/f21\_ncmapss\_cert\_floor.py}, \code{prereg-v24}):

\begin{center}\small
\begin{tabular}{@{}lrr@{}}
  \toprule
  Life fraction & irreducible share, N-CMAPSS & irreducible share, ours \\
  \midrule
  $50\,\%$ & $0.68$ & $0.87$ \\
  $70\,\%$ & $0.72$ & $0.85$ \\
  $90\,\%$ & $0.82$ & $0.88$ \\
  \bottomrule
\end{tabular}
\end{center}

Absolute floors are not comparable --- N-CMAPSS units live tens of cycles where ours live
thousands --- but the \emph{share} is, and in both fleets most of the remaining-life spread
between engines in identical true condition is irreducible. \textbf{F11's central claim
survives on data this project did not generate.} The trend differs in detail: ours is roughly
flat across life, N-CMAPSS rises, which is worth noting and not worth explaining away here.

\subsection{The certificate does not port --- and the control is the reason}

F10 needs three things N-CMAPSS does not ship: an influence matrix (estimated in
\cref{sec:f20-ncmapss}), a noise covariance (estimated from within-flight lag-1 residuals), and
an estimator to be honest \emph{about} (a WLS gas-path solve). Run that way, the certificate
appears to work beautifully: Spearman $\rho = 0.842$ ($p = 0.0022$) between the bound and the
estimator's true per-direction error, better than the $0.70$ of the original.

It is an artifact, and one control shows it. The bound and the estimator are computed from the
\emph{same} matrix. Column-shuffling that matrix destroys the physics while keeping the coupling
and the multiset of column norms --- and the shuffled control reproduces the correlation almost
exactly, $\rho = 0.830$ against the real $0.842$.

\keyidea{Both the CRB for direction $j$ and the estimator's error in direction $j$ grow when column $j$ of $\mathbf{H}$ is weak, \emph{whether or not that column is the physically right one}. A ranking correlation between them can therefore be produced by a matrix with no physics in it. \textbf{The N-CMAPSS port measures the coupling, not the certificate, and is withdrawn.}}

\subsection{Running the same control on our own F10}

That mechanism is not specific to N-CMAPSS. In \cref{ch:breakthrough} the certificate is built
from the pyCycle influence matrix and the Kalman estimator uses the same matrix, so the same
coupling exists, and \textbf{F10 never had this control}. Contribution C8 is the project's
most-promoted result; the control is cheap; it had to be run
(\code{scripts/f22\_f10\_shuffle\_control.py}, \code{prereg-v25}).

\begin{center}\small
\begin{tabular}{@{}lrrr@{}}
  \toprule
  & real $\rho$ & control median & empirical $p$ \\
  \midrule
  this study (pyCycle CFM56-7B) & $\mathbf{0.697}$ & $0.242$ & $0.085$ \\
  N-CMAPSS port                 & $0.842$ & $0.830$ & --- \\
  \bottomrule
\end{tabular}
\end{center}

\begin{figure}[htbp]
  \centering
  \includegraphics[width=0.64\textwidth]{shuffle_control}
  \caption{The honesty correlation against its own null. Grey spans the 5th--95th percentile of
  200 column-shuffled draws, in which the influence matrix carries no physics. Auto-generated
  (\protect\code{prereg-v25}).}
  \label{fig:shuffle-control}
\end{figure}

\textbf{The reimplementation is faithful} --- $\rho = 0.697$ against the published $0.70$ --- so
what follows is about the original claim, not about a coding difference.

\textbf{The verdict is genuinely mixed, and both halves matter.} The control median is
$\mathbf{0.242}$, not zero: a physics-free matrix already ranks somewhat, so \emph{part} of the
published $0.70$ was always the coupling, and that context was never in this document. But the
real value sits far above that median, so the certificate does carry real signal. What it does
not do is clear significance: the null distribution has a long tail whose 95th percentile
reaches $0.722$, \emph{above} the real value, giving $p = 0.085$.

\keyidea{H10.1 is neither confirmed nor refuted by this control --- it is \textbf{under-powered}. With ten health directions, a rank correlation cannot separate a certificate that works from a coupling that mimics one. The honest status of $\rho = 0.70$ is: real signal, partly attributable to a shared matrix, and never significant against its own null.}

\paragraph{What follows for C8.} The claim is qualified rather than withdrawn, and the
qualification travels with the number: every place this document quotes $\rho = 0.70$ now
quotes the null beside it (\cref{sec:contributions,ch:breakthrough,sec:questions-answered}).
The contrast with N-CMAPSS is the useful part --- there the certificate added \emph{nothing}
over the coupling, here it adds a great deal, which is evidence that a well-founded influence
matrix matters and an estimated one is not a substitute.

\paragraph{Why this was not caught earlier, stated plainly.} Every later line in this chapter
carries a control arm --- \cref{sec:f17-hybrid} has one that fired, \cref{sec:f19-cert} has one
that inverted a post-hoc, \cref{sec:f15-instrument} has a fraud check that validated a null.
F10 predates that discipline and never had one. The lesson is not about this certificate: it is
that \textbf{the control is the part of an experiment worth designing first}, and the results
here that survived scrutiny are exactly the ones that had one.

\subsection{Breaking the coupling by design}
\label{sec:f23-decoupled}

The shuffle test is blunt because shuffling $\mathbf{H}$ moves the bound \emph{and} the
estimator. The clean version removes the shared object instead: replace the Kalman with a
\textbf{learned} estimator --- a sequence model mapping the deviation trajectory to the ten
health parameters, trained against ground truth --- which never sees $\mathbf{H}$ at all
(\code{scripts/f23\_decoupled\_certificate.py}, \code{prereg-v26}, ten seeds).

Now permuting the certified precisions changes the claim under test and nothing else, so the
permutation is an honest null. It is also nearly free: nothing is retrained, so ten thousand
permutations cost less than one of the two hundred shuffles the previous test needed.

\begin{center}\small
\begin{tabular}{@{}lrrrl@{}}
  \toprule
  Estimator & $\rho$ & null median & null p95 & $p$ \\
  \midrule
  Kalman (shares $\mathbf{H}$ with the bound) & $0.697$ & $0.242$ & $0.722$ & $0.085$ \\
  learned (never sees $\mathbf{H}$)           & $0.455$ & $\mathbf{0.006}$ & $0.552$ & $0.100$ \\
  \bottomrule
\end{tabular}
\end{center}

\begin{figure}[htbp]
  \centering
  \includegraphics[width=0.86\textwidth]{decoupled}
  \caption{Left: certified precision against the learned estimator's true per-direction error.
  Right: each test against its own null --- square is the null, circle the measured value.
  Removing the shared matrix collapses the null from $0.242$ to $0.006$. Auto-generated
  (\protect\code{prereg-v26}).}
  \label{fig:decoupled}
\end{figure}

\textbf{The design worked, and the claim still cannot be proven.} The null collapses from
$0.242$ to $\mathbf{0.006}$ --- essentially zero, as an honest permutation should be --- so the
coupling really is gone. The certificate still ranks the error at $\rho = 0.455$, consistently
across all ten seeds ($0.18$ to $0.56$, every one positive), so this is not one lucky model.
But $p = 0.100$: with ten directions, $\rho$ would have to reach about $0.55$ to clear
significance.

\keyidea{The two experiments are consistent, and together they decompose the original number: of F10's $\rho = 0.70$, roughly $0.24$ was the shared matrix and roughly $0.46$ is signal that survives when the coupling is removed. What blocks the claim is not the certificate --- it is that \textbf{ten health directions is too few for a rank test}, and ten is fixed by the physics, not by how much fleet data one collects.}

\paragraph{The honest status of C8, stated once and for all.} The certificate is not shown to
be dishonest. It is shown to be \textbf{unprovable by the test F10 chose}. That is a defect in
the validation design, not in the instrument, and it is a more useful finding than the earlier
ambiguity: before this section one could not distinguish ``the certificate does not work'' from
``the certificate cannot be measured this way''. Now one can.

\subsection{Fixing the scale, and the second wall behind it}
\label{sec:f24-scale}

The previous subsection named the way out: stop ranking, measure magnitude. That needs the
CRB's absolute scale to be trustworthy, which \cref{sec:bt-honesty} shows it is not, so the
order is calibrate then test (\code{scripts/f24\_crb\_scale\_magnitude.py},
\code{prereg-v27}). Errors still come from the learned estimator of
\cref{sec:f23-decoupled}, so the decoupling is preserved.

\paragraph{Part 1 --- the scale, and it is not a scale.} Fitting
$\log|\mathrm{err}| = a + b \log \mathrm{CRB}$ on \emph{validation} engines only:

\begin{center}\small
$\log|\mathrm{err}| = -2.188 + 0.382\,\log \mathrm{CRB}$
\end{center}

Two readings, and the second is the important one. The bound understates achievable error by
$\mathbf{8.9\times}$ at $\mathrm{CRB}=1\,\%$ --- expected, since a Cram\'er--Rao bound is a floor
for \emph{unbiased} estimators. But the exponent is $\mathbf{0.382}$, not $1$. \textbf{The
distortion is not a scale error, it is a shape error in magnitude}: a single multiplicative
constant cannot repair the bound, because the relationship between certified precision and
achievable error is not proportional. After a two-parameter calibration, $90\,\%$ of test
errors still fall only within a factor of $\mathbf{5.6}$ of prediction --- a precision forecast
with six-fold uncertainty decides nothing operationally.

\paragraph{Part 2 --- the powered statistic, and the gate that disqualifies it.} The certificate
is per engine: it accumulates Fisher information over that engine's own flown conditions, so
engines that flew differently should earn different bounds. Differencing out each direction then
gives a test with $n = 200$ instead of $10$, immune to the ten-direction ceiling.

That test was gated in advance on whether the CRB varies between engines at all, and
\textbf{the gate fails}: the within-direction coefficient of variation of the CRB is
$\mathbf{1.15\,\%}$ against a $5\,\%$ threshold.

\begin{center}\small
\begin{tabular}{@{}lrrl@{}}
  \toprule
  Statistic & value & $p$ & $n$ \\
  \midrule
  between direction (F10's) & $\rho = 0.370$ & $0.293$ & $10$ \\
  within direction & $r = 0.145$ & $0.041$ & $200$ \\
  within direction & $\rho = 0.184$ & $0.009$ & $200$ \\
  \bottomrule
\end{tabular}
\end{center}

\keyidea{The within-direction test returns $p = 0.041$ and $p = 0.009$, and \textbf{those numbers are not reported as evidence}, because the pre-declared gate says they cannot be. With a predictor that varies by one per cent across engines, a correlation of $0.145$ at $n = 200$ is what a large sample manufactures, not what a real effect looks like. Quoting them would be the same error the shuffled-matrix control exists to catch, wearing different clothes.}

\paragraph{Why the fleet cannot answer this.} The gate fails for a reason that is a property of
SynCFM56 rather than of the certificate: every engine flies a near-identical N1 profile, so the
accumulated Fisher information comes out nearly the same for all of them. The certificate is
per-engine in principle and near-constant in practice, and with no variation in the predictor
there is nothing to regress against.

\paragraph{What would settle C8, concretely.} A fleet with \textbf{deliberate mission
diversity} --- engines flying materially different route mixes, so that different histories earn
genuinely different bounds. Then the within-direction test has a predictor that moves, the
ten-direction ceiling is irrelevant, and certified precision can be regressed against achieved
error with real power. The lever is known to exist: \cref{ch:tomography} already showed that
changing \emph{one} report condition moves separability by $79\,\%$, so operating-history
diversity has first-order leverage on this geometry. That is a generator change plus a re-run,
which is a study of its own and is recorded as the one experiment that would resolve the
project's most-promoted claim.

\keyidea{After three tests, C8's status is precise rather than vague. The certificate is not shown to be dishonest. It is shown to be unmeasurable \emph{on this fleet}: the ranking statistic is capped at ten directions by the physics, and the magnitude statistic is disabled by a fleet whose engines all fly the same way. The fix is a more diverse fleet, not a better estimator.}

\section{L-ICM: the sensitivity study the premise owed}
\label{sec:f-icm}

\paragraph{Why this one is different.} Every line above tests a \emph{result}. This one tests
a \emph{premise} --- and the load-bearing one. \Cref{sec:limitations} closes the model chapter
with an assertion: the declared limitations (generic maps, no transients, approximate speed
labels) do not endanger the comparison, ``which rests on the model's sensitivity structure
(the ICM), not on its absolute values.'' Everything downstream leans on it: the H2 wall is a
statement about ICM geometry, the confusable subset is defined by ICM angles, the F10
certificate inverts the same matrix. Yet the assertion was never measured. If a calibration
error inside the contract's own tolerance moved the signature angles appreciably, the
observability verdict would inherit the twin's uncertainty and every claim built on it would
need that caveat. This line performs the classical GPA sensitivity study on our own premise
(\code{docs/prereg-v13.md}, tag \code{prereg-v13}).

\paragraph{What is perturbed.} Each calibration knob is moved inside the tolerance its own
contract declares (\cref{tab:tcds}, \cref{ch:decisions}), one at a time and then jointly at
adverse signs: the HPC design pressure ratio and the HP map speed reference ($\pm3$\,\%, the
two quantities \cref{sec:calibration} actually calibrated), the design thrust and rating
temperature ($\pm3$\,\%, the class-[A] and class-[D] values), the design bypass ratio
($\pm0.2$, the EEDB tolerance) and all five component efficiencies ($\pm2$\,\% relative ---
the generic-map limitation made numeric). The cockpit ICM is recomputed from scratch at
cruise and at hot-day takeoff for every perturbation, and the observability verdict is
re-read: rank, minimum signature angle, the identity and membership of the confusable set,
and the angle by which each individual fault signature rotates. The same run sweeps the
finite-difference step ($\pm0.25$ / $\pm0.5$ / $\pm1.0$\,\%), until now a reasoned
convention rather than a measured one.

\begin{table}[htbp]
  \centering\small
  \caption{L-ICM pre-registered verdicts (\texttt{prereg-v13}): does the observability
  verdict survive calibration error? Auto-generated.}
  \label{tab:icm-robustness}
  % AUTOGENERATED by scripts/make_report_assets.py — do NOT edit by hand
\begin{tabular}{@{}>{\raggedright\arraybackslash}p{0.09\textwidth}>{\raggedright\arraybackslash}p{0.42\textwidth}>{\raggedright\arraybackslash}p{0.43\textwidth}@{}}
\toprule
ID & Frozen criterion & Verdict \\
\midrule
H-ICM.1 & Rank, minimum angle and most-confusable pair unchanged under every calibration perturbation & \textbf{confirmed} --- min.\ angle stays in 0.94--1.45$^\circ$, rank 3, top pair unchanged \\
H-ICM.2 & Median signature-direction shift $\leq 5^\circ$ and the confusable set unchanged in $\geq 90$\,\% of perturbations & \textbf{refuted} --- median 0.61$^\circ$, worst 31.69$^\circ$; set unchanged in 81\,\% \\
H-ICM.3 & The $\pm0.5$\,\% finite-difference step is in the linear regime ($<2^\circ$, $<5$\,\% vs 0.25/1.0\,\%) & \textbf{refuted} --- 5.99$^\circ$, 4.22\,\% magnitude \\
\bottomrule
\end{tabular}

\end{table}

\begin{figure}[htbp]
  \centering
  \includegraphics[width=\textwidth]{icm_robustness}
  \caption{L-ICM. Left: the minimum cockpit signature angle under every calibration
  perturbation, at both snapshot conditions; dotted lines are the nominal values, the dashed
  line the $15^{\circ}$ confusability threshold. Right: how far each fault signature rotates
  (dot = median over the ten health parameters, tick = worst one) against the pre-registered
  $5^{\circ}$ bound. Auto-generated.}
  \label{fig:icm-robustness}
\end{figure}

\paragraph{The premise holds --- H-ICM.1 confirmed.} Under all twenty-six
perturbation-condition combinations the cockpit matrix keeps rank 3, the minimum signature
angle stays between $0.94$ and $1.45^{\circ}$ (nominal $1.32^{\circ}$ at cruise,
$1.25^{\circ}$ at hot-day takeoff), and the most-confusable pair remains
$\eta_{\mathrm{HPC}}\!\sim\!\eta_{\mathrm{HPT}}$ every single time --- including under the
joint adverse draw that moves every knob at once. The H2 wall is not an artifact of how the
twin was calibrated. Neither is it fragile to \emph{how} the matrix is computed: the finite
differences are insensitive to the design thrust and the map speed reference almost entirely
(median shifts of $0.00$--$0.03^{\circ}$ --- the speed reference is a labelling convention, so
the physics behind the sensitivities never notices it), and even the rating temperature, the
most influential knob, leaves the verdict intact.

\paragraph{But the margins are not --- H-ICM.2 refuted, and worth the refutation.} The median
signature rotation is $0.61^{\circ}$, an order of magnitude inside the pre-registered
$5^{\circ}$ bound, and no perturbation exceeds $1.86^{\circ}$ at the median. The clause that
fails is the second one: the confusable set is unchanged in only $81\,\%$ of perturbations,
under the frozen $90\,\%$ requirement. Five perturbations at hot-day takeoff move a pair
across the $15^{\circ}$ line --- $\Gamma_{\mathrm{LPC}}\!\sim\!\Gamma_{\mathrm{LPT}}$ enters,
$\eta_{\mathrm{fan}}\!\sim\!\eta_{\mathrm{LPC}}$ leaves --- taking the count from seven to six
or eight. The mechanism is the same one that produces the $31.7^{\circ}$ worst-case rotation,
and it has a name in this document already: $\Gamma_{\mathrm{LPC}}$ carries one of the
weakest cockpit signatures (a column norm of $0.27$ against fan flow capacity's $2.68$), so
its \emph{direction} is numerically soft and swings under any change to the cycle. The
directions whose geometry is unstable are precisely the directions the F10 certificate
already labels unobservable, and whose diagnosis \cref{ch:gpa-practice} already showed to be
worthless. The refutation therefore sharpens rather than damages the picture: \textbf{the
load-bearing verdict is calibration-invariant, the marginal membership of the confusable list
is not}, and any future work that treats the seven-pair count as a fixed property of the
engine --- rather than the core pair plus a soft fringe --- would be over-reading it.

\paragraph{And the step needs saying out loud --- H-ICM.3 refuted.} Recomputing the matrix at
$\pm0.25\,\%$ and $\pm1.0\,\%$ moves column magnitudes by at most $4.2\,\%$ (inside the frozen
$5\,\%$) but rotates the worst column by $5.99^{\circ}$, three times the frozen $2^{\circ}$.
Once more the offender is $\Gamma_{\mathrm{LPC}}$, with $\Gamma_{\mathrm{LPT}}$ second at
$2.15^{\circ}$; for the seven parameters that matter operationally the shift is a small
fraction of a degree. The honest statement is therefore narrower than the one the register
carried before this line ran: the $\pm0.5\,\%$ central difference is in the linear regime
\emph{for the observable directions}, and is not a well-posed quantity for the two near-null
ones --- because a near-zero vector has no stable direction to measure. That is a property of
the engine, not of the step.

\paragraph{What it costs, and what it buys.} Two of three frozen criteria failed, and both
failures trace to one cause that the study had no way to anticipate --- weak columns have
unstable directions --- which is exactly the kind of thing pre-registration is supposed to
catch rather than let an author quietly discover and then present as if it had been expected.
The premise the chapter set out to test survives in the form that matters: every claim in this
document that rests on \emph{the confusable core} (H2, the certificate, the sensor
recommendation) is calibration-robust, and the claims that rest on the exact confusable
\emph{count} are now known not to be.

\section{Closing C6: the noise axis of the determinability map}
\label{sec:noise-sweep}

\paragraph{The axis that was missing.} Contribution C6 promised a determinability map over
three axes --- \emph{sensors} $\times$ \emph{noise} $\times$ \emph{data size} --- telling an
operator when each approach wins. Two were delivered and one was not. The sensor axis runs
through the whole document (cockpit versus extended, from \cref{sec:icm} to
\cref{ch:gpa-practice} to the certificate); the data-size axis is the H4 ablation at
$10/25/100\,\%$ of the training engines (\cref{sec:ai-hybrid}). But every result in this
document, without exception, lives at \emph{one} sensor-quality setting: the literature
values frozen in the fault catalog. Whether the verdicts survive better or worse
instrumentation was, until this section, unknown --- and unknown in the direction that
matters, because sensor quality is the one axis an operator can buy along.

\paragraph{Design.} The fleet is regenerated with every sensor $\sigma$ scaled by
$0.5$, $1$, $2$ and $4$, holding everything else byte-identical: same seed, same health
trajectories, same splits, same fault magnitudes and onsets. Measurement quality is the only
variable that moves, so any change in a score is attributable to it alone; the $1\times$ level
reproduces the frozen v1.1 fleet exactly (median life $15\,139$ cycles either way). Both
families are scored at their \textbf{F5-tuned configurations, unchanged}. That choice is
disclosed rather than buried: the map therefore measures how the systems an operator would
have \emph{fielded} degrade as sensors change, not the best score reachable at each level ---
which would need eight further Optuna campaigns and answers a different question.
Pre-registered as \code{docs/prereg-v14.md}, tag \code{prereg-v14}.

\begin{table}[htbp]
  \centering\small
  \caption{The determinability map, noise axis: both tuned families across a
  $0.5$--$4\times$ range of sensor noise. Auto-generated.}
  \label{tab:noise-sweep}
  % AUTOGENERATED by scripts/make_report_assets.py — do NOT edit by hand
\begin{tabular}{@{}lrrrrrr@{}}
\toprule
& \multicolumn{2}{c}{RUL RMSE @90\,\% [cy]} & \multicolumn{2}{c}{confusable isolation} & \multicolumn{2}{c}{detection recall} \\
\cmidrule(lr){2-3}\cmidrule(lr){4-5}\cmidrule(lr){6-7}
Sensor noise & trad. & AI & trad. & AI & trad. & AI \\
\midrule
$0.5\times$ & 1955 & 867 & 0.38 & 0.31 & 0.48 & 0.52 \\
$1\times$ (as fielded) & 1981 & 862 & 0.31 & 0.31 & 0.13 & 0.48 \\
$2\times$ & 2184 & 860 & 0.23 & 0.46 & 0.04 & 0.39 \\
$4\times$ & 2786 & 815 & 0.15 & 0.46 & 0.00 & 0.35 \\
\bottomrule
\end{tabular}

\end{table}

\begin{table}[htbp]
  \centering\small
  \caption{Noise-axis pre-registered verdicts (\texttt{prereg-v14}). Auto-generated.}
  \label{tab:noise-verdicts}
  % AUTOGENERATED by scripts/make_report_assets.py — do NOT edit by hand
\begin{tabular}{@{}>{\raggedright\arraybackslash}p{0.08\textwidth}>{\raggedright\arraybackslash}p{0.46\textwidth}>{\raggedright\arraybackslash}p{0.40\textwidth}@{}}
\toprule
ID & Frozen criterion & Verdict \\
\midrule
H-N.1 & The AI prognosis advantage holds at every noise level & \textbf{confirmed} --- AI ahead at 4 of 4 levels \\
H-N.2 & Confusable-isolation gap between families $\leq 10$\,pp everywhere & \textbf{refuted} --- gap 0--31\,pp \\
H-N.3 & Halving the noise does not lift either family past $50$\,\% confusable isolation & \textbf{confirmed} --- traditional 0.38, AI 0.31 \\
\bottomrule
\end{tabular}

\end{table}

\begin{figure}[htbp]
  \centering
  \includegraphics[width=\textwidth]{noise_sweep}
  \caption{How the two tuned families degrade as sensor noise is scaled from half to four
  times its nominal value, on prognosis (left), confusable isolation (centre) and detection
  (right). The $1\times$ point is the fleet every other chapter uses. Auto-generated.}
  \label{fig:noise-sweep}
\end{figure}

\paragraph{Prognosis is noise-proof, and the advantage widens --- H-N.1 confirmed.} The AI's
RUL error at $90\,\%$ of life is essentially \emph{flat} across a factor of eight in sensor
noise: $867$, $862$, $860$, $815$ cycles, a spread smaller than the seed-to-seed variation of
the training itself ($\pm44$--$66$ cycles). The traditional margin extrapolation, over the
same range, degrades from $1\,955$ to $2\,786$. The advantage therefore does not merely
survive worse sensors --- it \textbf{grows}, from $2.3\times$ at half noise to $3.4\times$ at
quadruple. The mechanism is the one H3 rested on all along: remaining life is a
temporal-aggregation task, and a sequence model that reads a thousand flights averages
per-flight noise into insignificance, while a Theil--Sen slope through a noisier margin series
gets both a worse slope and a worse extrapolation from it. For an operator, this is the
robustness statement that matters commercially: the prognosis business case of
\cref{ch:economics} does not depend on sensor quality.

\paragraph{Detection tells the same story more sharply.} Traditional recall collapses from
$0.48$ at half noise to $0.13$ at nominal, $0.04$ at double and $0.00$ at quadruple --- with
false alarms falling to zero alongside it, because the tuned persistence rules simply stop
firing at all. The AI detector goes $0.52 \to 0.48 \to 0.39 \to 0.35$ at a constant single
false alarm. This puts the H1 result (\cref{sec:res-h1}) in a new light: the tuned traditional
detector is not merely worse at nominal noise, it is \emph{brittle} to sensor quality, and it
happened to be measured just past the edge of where its thresholds work. Its thresholds were
selected on a validation split with two clean engines, and that fragility is now visible as a
slope rather than as a single unlucky number.

\paragraph{The tie does not hold everywhere --- H-N.2 refuted, and carefully.} The
confusable-isolation gap between the families is $7.7$\,pp at half noise, exactly $0.0$\,pp at
nominal (the H2 tie, reproduced), then $23.1$ and $30.8$\,pp at double and quadruple noise ---
the AI ahead. Beyond $10$\,pp twice, so the frozen criterion fails and stays failed. Two things
must be said about it, in this order. First the honest caveat: there are \textbf{thirteen}
confusable test episodes, so $10$\,pp is $1.3$ episodes and the $23$\,pp gap is three of them.
A threshold that fine was mis-calibrated against the sample it would be judged on --- the same
error that retired H7.3 --- and no claim of statistical separation is made here. Second, the
direction is mechanistically what one would expect and is worth recording as a hypothesis for
someone with more episodes: the nearest-signature rule decides from a single measured step, so
noise hits it undiluted, while the learner's features average step, slope and level over
hundreds of cycles. \emph{If} that holds up, the practical corollary is that the physics rule's
edge at nominal noise is conditional on good instrumentation, and degrades where the learner
does not.

\paragraph{And the wall is geometric, not noise-driven --- H-N.3 confirmed.} This is the
result C6 was owed. Halving every sensor $\sigma$ --- a genuinely expensive instrumentation
upgrade --- leaves confusable isolation at $0.38$ (traditional) and $0.31$ (AI), both below
the coin-flip line and both far from the $0.92$ that \emph{different} sensors deliver
(\cref{sec:f8-lh2}). Quieter sensors of the same set do not break the wall, because two
signatures $1.3^{\circ}$ apart stay $1.3^{\circ}$ apart however precisely they are measured:
noise sets how far apart two faults must be to be separable, but the geometry decides whether
they ever are. The purchase order the document has been writing since \cref{sec:res-h2} is
therefore specific in a way it could not be before: \textbf{buy probes at new stations, not
better probes at the old ones}.

\paragraph{The map, assembled.} With this axis, C6 is complete and reads as three different
answers to ``when does each approach win''. Along \emph{data size}, the learner saturates at
about seventeen engines and physics features never help (\cref{sec:ai-hybrid,sec:f8-l6}).
Along \emph{sensors}, both families are bounded by the same geometry, and the extended set
moves both (\cref{ch:gpa-practice,sec:f8-lh2}). Along \emph{noise}, the two families diverge:
learning degrades gracefully and classical thresholds fall off a cliff. The single sentence an
adopting organization should take away is that \emph{sensor quality changes who wins, sensor
placement changes what is knowable, and only the second one has a ceiling}.

\section{The F8 program so far}
\label{sec:f8-summary}

\Cref{tab:f8-summary} collects the ten executed lines and their verdicts. The shape of the
table is the point: capability-extending lines (L1, L2, L4, L5) succeed, while lines that
probe the confusable-isolation wall (L7, L9) return partial results that reinforce it, the
one physics-injection line (L6) is a clean negative, and the one line that audits the
project's own premise (L-ICM) confirms it where it is load-bearing and refutes two of its
own frozen margins. The honest ledger is what a research program looks like when its gates
are allowed to fail.

\begin{table}[htbp]
  \centering\small
  \caption{The ten executed F8 limitation-overcoming lines and their pre-registered verdicts.}
  \label{tab:f8-summary}
  \begin{tabular}{>{\raggedright\arraybackslash}p{0.085\textwidth}>{\raggedright\arraybackslash}p{0.40\textwidth}>{\raggedright\arraybackslash}p{0.435\textwidth}}
    \toprule
    Line & Question & Verdict \\
    \midrule
    L1 & A differentiable twin at pyCycle fidelity? & \textbf{Yes} --- 13--19$\times$ closer than the linearization (\cref{sec:f8-l1}) \\
    L2 & A nonlinear fleet from it? & \textbf{Yes} --- SynCFM56 v2, 3--7$\times$ more faithful, gates re-passed (\cref{sec:f8-l2}) \\
    L5 & Is the RUL win architecture-robust? & \textbf{Yes} --- GRU, TCN, Transformer all beat traditional (\cref{sec:f8-l5}) \\
    L4 & Is the recoverable margin fraction estimable? & \textbf{Yes} --- $R^2 = 0.86$ from trajectory shape (\cref{sec:f8-l4}) \\
    L6 & Does a twin-residual feature help RUL (H4)? & \textbf{No} --- refuted a second, independent way (\cref{sec:f8-l6}) \\
    L7 & Can an augmented Kalman fix a drifting sensor? & \textbf{Partly} --- tracks the drift, cannot un-corrupt the cockpit diagnosis (\cref{sec:f8-l7}) \\
    L9 & Can the PCS metric be validated? & \textbf{Partly} --- tracks correctness weakly; no strongly-physical classifier to test against (\cref{sec:f8-l9}) \\
    L-H2 & Does breaking the wall need a \emph{real} sensor? & \textbf{Yes} --- real station sensors lift confusable isolation 0.15$\to$0.92; a virtual one adds nothing (\cref{sec:f8-lh2}) \\
    L-RUL & Does an advanced classical prognostic narrow H3? & \textbf{Yes, but} --- it halves the traditional error (1981$\to$1118); the AI still wins 858 (\cref{sec:f8-lrul}) \\
    L-RUL/sweep & Does the AI beat it at \emph{every} life fraction? & \textbf{No} --- only at 90\,\% ($1.34\times$, $p=0.006$); a tie at 70\,\% ($1.01\times$). The learning advantage is late-life (\cref{sec:f8-lrul-sweep}) \\
    L-UQ & Is the 5.9$\times$ interval advantage the model or the calibration? & \textbf{Mostly calibration} --- at equal coverage the AI leads $1.34\times$ over advanced classical, $2.3\times$ over practice (\cref{sec:f-uq}) \\
    L-MECH & Is the wall instantaneous? Can a history attribute mechanism? & \textbf{Yes} --- 3 of 5 mechanisms recover from trajectory shape, \code{hot\_section} at $R^2\,0.78$; but learning beats hand features on only 2 of 5 (\cref{sec:f13-mech}) \\
    L-INST & Can either family tell a drifting sensor from a real fault? & \textbf{No} --- both at chance (AUC $0.61$ / $0.46$, CIs spanning $0.5$); tracking a bias is not attributing one (\cref{sec:f15-instrument}) \\
    L-EXT & Does a published deep-learning PHM result survive these standards? & \textbf{Partly} --- it replicates ($14.56$ vs $14.12$) but fails to train on 5 of 10 runs, and its 15-way ranking is inside seed noise (\cref{sec:f14-ext}) \\
    L-HYB & Does the certificate work as a physics-injection channel? & \textbf{Refuted by its control} --- both physics arms beat pure ($p<0.001$), so the gain is channels. But injection \emph{works} on attribution where H4/L6 refuted it on prognosis (\cref{sec:f17-hybrid}) \\
    L-CERT & Does gating the GPA estimate with the certificate help? & \textbf{No --- it hurts} ($-0.119$, $t=-3.72$) at equal channel count; the CRB as input adds nothing ($+0.007$). L-HYB's post-hoc was the channel count (\cref{sec:f19-cert}) \\
    L-BIDIR & Is the published Bi-LSTM's edge direction or capacity? & \textbf{Capacity, on a GRU cell} --- 94\,cy of RUL error at doubled parameters ($p=0.016$), nothing at equal. \textbf{Does not transfer to LSTM} (sign reverses), and that comparison is confounded by GRU-tuned hyperparameters (\cref{sec:f18-bidir,sec:f18-lstm}) \\
    L-EXT2 & Is the confusable wall an artifact of our own deck? & \textbf{No} --- on NASA's C-MAPSS, a different engine and code, the pair sits at $0.66^{\circ}$ on cockpit-class channels and opens to $20.9^{\circ}$ with full instrumentation; ours are $1.3^{\circ}$/$26.7^{\circ}$ (\cref{sec:f20-ncmapss}) \\
    L-DECOUP & Is the certificate broken, or unmeasurable this way? & \textbf{Unmeasurable} --- with a learned estimator that never sees $\mathbf{H}$ the null collapses to $0.006$ and $\rho=0.455$ holds on all ten seeds, but $p=0.100$: ten directions is too few for a rank test (\cref{sec:f23-decoupled}) \\
    L-SCALE & Can magnitude be tested instead of ranking? & \textbf{Not on this fleet} --- the CRB understates error $8.9\times$ with exponent $0.382$ (a shape error, not a scale one), and the powered test is disqualified by its own gate: the CRB varies only $1.15\,\%$ between engines (\cref{sec:f24-scale}) \\
    L-CERT2 & Do the certificate and the floor port to external truth? & \textbf{Floor yes, certificate no} --- irreducible share $0.68$--$0.82$ vs ours $0.87$--$0.88$; the certificate's $\rho=0.842$ on N-CMAPSS is matched by a physics-free null at $0.830$, and the same control on \emph{our} F10 gives $\rho=0.697$ against a null median $0.242$ at $p=0.085$ (\cref{sec:f21-port}) \\
    Gate T & Would transients open the confusable angle? & \textbf{Not via operating points} --- the whole envelope gives $1.38^{\circ}$, below cruise alone; true dynamics are priced, not ruled out (\cref{sec:gate-t}) \\
    L-ICM & Does the observability verdict survive calibration error? & \textbf{Yes at the core, no at the fringe} --- the confusable \emph{pair} is invariant, the confusable \emph{count} is not (\cref{sec:f-icm}) \\
    C6/noise & Do the verdicts survive better or worse sensors? & \textbf{Prognosis yes, detection emphatically, the tie no} --- and halving the noise does not break the wall (\cref{sec:noise-sweep}) \\
    \bottomrule
  \end{tabular}
\end{table}
